\chapter{GraphQL: Schema-Driven APIs}
\label{ch:graphql}

\section*{Executive Summary}
\addcontentsline{toc}{section}{Executive Summary}

REST asks the server to guess what clients need; GraphQL makes the client
say it. That single inversion --- the client composes the response shape,
the server guarantees it against a typed schema --- is the whole idea, and
everything else in this chapter is a consequence of it. The consequences
are larger than they look. A schema that clients query against directly is
a \emph{contract with an algebra}: fields compose, fragments factor out
shared selection, nullability propagates errors along formally specified
paths, and the type system is introspectable at runtime. The same property
that makes GraphQL delightful for a mobile team --- one endpoint, one
round trip, exactly the fields requested --- makes it dangerous by
default: the client, not the server, decides how much work a request
costs. A REST endpoint's worst case is bounded by its implementation; a
GraphQL endpoint's worst case is bounded by the attacker's imagination,
unless you bound it yourself.

This chapter treats GraphQL as a specification first and a library second.
We work through the October 2021 specification's type system \cite{graphqlspec}
--- scalars, objects, interfaces, unions, enums, input objects, and the
list/nullability algebra --- and formalize the rule that makes GraphQL
error handling unlike anything in REST: when a non-null field errors, the
\texttt{null} bubbles upward to the nearest nullable position, and the
response arrives as \emph{partial data} in an HTTP 200 with a populated
\texttt{errors} array. We trace an operation through the parse, validate,
and execute phases; we dissect the request/response wire shape; and we
confront the two production taxes every GraphQL server pays. The first tax
is the $N{+}1$ query problem, which field-level resolution creates by
construction: we reproduce it with an instrumented repository and watch a
single innocent query execute 11 store reads, then collapse it to 3 with
per-request DataLoader batching --- measured, not asserted. The second tax
is cost control: depth limiting, complexity analysis, operation
allow-lists, and Automatic Persisted Queries, which we implement as a
pre-execution gate that rejects expensive documents before a single
resolver runs.

The reference implementation is a Strawberry-GraphQL catalog API for the
fictional Northwind Robotics parts universe: Relay-style cursor
connections, typed mutation errors via union payloads, field-level error
extensions, an instrumented repository, a cost-control layer, an
APQ-capable httpx client, and a 21-test in-process suite --- all
deterministic, all offline, all runnable on a Windows workstation with
PowerShell. The chapter closes with an honest GraphQL-versus-REST-versus-gRPC
decision matrix and a scoped primer on federation, because the correct
answer to ``should we use GraphQL?'' is sometimes ``no,'' and a principal
engineer should be able to say when.

\begin{keyconceptbox}
A GraphQL schema is a \emph{queryable contract}: clients compose responses
from typed fields, and the server guarantees shape, nullability, and error
semantics per the specification \cite{graphqlspec}. The power inversion is
real --- the client chooses the work --- so production GraphQL is defined
less by resolvers than by its guardrails: depth and complexity budgets,
persisted-query allow-lists, gated introspection, per-request DataLoaders,
and an error model that treats partial success as a first-class outcome.
\end{keyconceptbox}

%% ============================================================================
\section{Learning Objectives \& Prerequisites}
\label{sec:ch12-objectives}

\subsection{Learning Objectives}

After completing this chapter and its lab, you will be able to:

\begin{enumerate}[label=\textbf{LO-\arabic*.}]
  \item Define every kind in the GraphQL type system --- scalar, object,
        interface, union, enum, input object, list, and non-null --- and
        predict the schema a Strawberry definition generates.
  \item State and apply the null-propagation rule formally: an error at a
        non-null position replaces that position with \texttt{null} and
        propagates to the nearest nullable ancestor, up to the
        \texttt{data} root.
  \item Explain the three operation types (query, mutation, subscription)
        and the parse $\to$ validate $\to$ execute pipeline, including
        where static cost analysis must hook in.
  \item Read and produce the GraphQL-over-HTTP wire shape: the
        \texttt{query}/\texttt{variables}/\texttt{operationName} request
        envelope and the \texttt{data}+\texttt{errors} response envelope,
        including partial responses.
  \item Reproduce the $N{+}1$ query problem with an instrumented
        repository, count the queries, and eliminate it with DataLoader
        batching --- including the one-DataLoader-per-request rule and its
        caching semantics.
  \item Implement Relay-style cursor connections
        (\texttt{edges}/\texttt{node}/\texttt{pageInfo},
        \texttt{first}/\texttt{after}) with opaque, decodable cursors.
  \item Distinguish field errors from top-level errors, use
        \texttt{extensions.code} as the machine-readable channel, and
        choose between union-typed mutation payloads and raised errors.
  \item Design a cost-control layer: depth limits, argument-weighted
        complexity estimation, operation allow-lists, Automatic Persisted
        Queries, and timeouts.
  \item Enumerate the GraphQL-specific attack surface --- introspection
        abuse, batching attacks, god-queries, resolver-side injection ---
        and map each to its mitigation.
  \item Choose between REST, GraphQL, and gRPC on evidence: cacheability,
        uploads, versioning philosophy, tooling gravity, and team shape.
\end{enumerate}

\subsection{Prerequisites}

This chapter assumes Chapters~0--8. Concretely:

\begin{itemize}
  \item \textbf{Chapter~0:} reproducible-measurement discipline and the
        synthetic Northwind Robotics universe.
  \item \textbf{Chapters~1--2:} HTTP request/response mechanics, POST
        semantics, and Python HTTP clients \cite{rfc9110,httpxdocs}.
  \item \textbf{Chapter~3:} REST resource modeling --- the baseline this
        chapter contrasts GraphQL against.
  \item \textbf{Chapter~4:} JSON serialization contracts; GraphQL response
        envelopes are exactly such contracts.
  \item \textbf{Chapter~5:} client-side credential and transport hygiene.
  \item \textbf{Chapter~6:} building FastAPI services; the GraphQL server
        here mounts on FastAPI \cite{fastapidocs}.
  \item \textbf{Chapter~7:} error modeling; GraphQL's
        \texttt{errors} array and \texttt{extensions} extend that chapter's
        problem-details thinking into a partial-success world.
  \item \textbf{Chapter~8:} cursor versus offset pagination; the Relay
        connection spec is that chapter's keyset idea in GraphQL clothing.
\end{itemize}

%% ============================================================================
\section{Deep Technical Mechanics \& Architectural Concepts}
\label{sec:ch12-mechanics}

\subsection{What GraphQL Actually Is}

Strip away the tooling and GraphQL is three documents. The \emph{schema}
is a typed, introspectable description of everything clients may ask for.
An \emph{operation} is a document the client sends, written in the same
type system, selecting exactly the fields it wants. The \emph{response}
is a JSON document whose shape mirrors the operation, key for key. The
server's job is to resolve each requested field against the schema,
nowhere else; the client's job is to tolerate partial data, because a
GraphQL response can be simultaneously a success and a failure.

\begin{definition}[GraphQL request execution]
Given a schema $S$, an operation document $O$, variable values $V$, and an
initial root value $R$, \emph{execution} is the function
$\mathrm{Execute}(S, O, V, R) \mapsto (\mathrm{data}, \mathrm{errors})$
where $\mathrm{data}$ mirrors $O$'s selection set, field errors are
collected into $\mathrm{errors}$ with their paths, and the null-propagation
rule (Definition~\ref{def:null-propagation}) governs how field errors
mutate $\mathrm{data}$ \cite{graphqlspec}.
\end{definition}

Three contrasts with REST matter architecturally. First, \textbf{one
endpoint, many shapes}: REST fixes the response shape per resource; GraphQL
fixes the \emph{universe} of shapes and lets the client project. Second,
\textbf{the type system is the documentation and the validator}: there is
no drift between docs and behavior, because the schema that documents the
API is the schema that validates the query. Third, \textbf{the client
controls cost}: a REST server author bounds each endpoint's work; a
GraphQL server author bounds the \emph{space of expressible queries}, which
is the subject of Section~\ref{sec:ch12-cost-control} and of half this
chapter's war stories.

\subsection{The Type System}
\label{sec:ch12-type-system}

GraphQL's type system has eight kinds, and production fluency means knowing
what each one \emph{refuses} to express.

\textbf{Scalars} are leaves: the built-ins \texttt{Int}, \texttt{Float},
\texttt{String}, \texttt{Boolean}, \texttt{ID}, plus custom scalars (a
\texttt{DateTime} is the canonical example) whose serialization the server
defines. \texttt{ID} serializes as a string and means ``opaque identifier;
do not parse.''

\textbf{Object types} are the workhorses: named sets of typed fields,
where field types may themselves be objects --- this recursion is what
makes a GraphQL API a \emph{graph}. \textbf{Interfaces} are abstract
objects: they declare fields that every implementing object must carry,
and selecting through an interface uses inline fragments for
concrete-type-only fields. \textbf{Unions} are interfaces without the
contract: member types share nothing, so \emph{every} field selection goes
through inline fragments --- which is exactly why unions are the right
tool for typed mutation results (Section~\ref{sec:ch12-errors}).
\textbf{Enums} are closed sets of string values, validated at the
operation layer. \textbf{Input objects} are objects that flow the other
way, into the server as arguments; they may contain scalars, enums, lists,
and other input objects, but never unions or interfaces --- input types
are closed under a smaller algebra than output types, by design.

\textbf{Lists and non-null} are type \emph{modifiers}, and their
composition is where the specification's subtlety lives. \texttt{[Review]}
means ``nullable list of nullable reviews''; \texttt{[Review!]!} means
``non-null list of non-null reviews.'' The four combinations are
semantically distinct, and choosing among them is the schema author's
error-propagation policy expressed as syntax.

\begin{table}[h]
  \centering
  \small
  \begin{tabularx}{\textwidth}{l l X}
    \toprule
    \textbf{SDL type} & \textbf{May be} & \textbf{Error in one element does} \\
    \midrule
    \texttt{[Review]}   & null, [null] & nulls only that element \\
    \texttt{[Review!]}  & null         & nulls the whole list \\
    \texttt{[Review]!}  & [null]       & nulls only that element; list never null \\
    \texttt{[Review!]!} & neither      & propagates past the list to the parent \\
    \bottomrule
  \end{tabularx}
  \caption{The four list/nullability combinations as error-containment policy.}
  \label{tab:ch12-nullability}
\end{table}

\begin{definition}[Null propagation on errors]
\label{def:null-propagation}
Let a field position $p$ with declared type $T$ raise an error during
execution. If $T$ is nullable, $p$ is set to \texttt{null} and an entry is
appended to \texttt{errors}. If $T$ is non-null, the error \emph{bubbles}:
$p$'s parent object is set to \texttt{null}, and the bubbling repeats at
each ancestor until a nullable position absorbs it; if no nullable
position exists, the entire \texttt{data} entry becomes \texttt{null}
\cite{graphqlspec}. Non-null is therefore not ``this cannot fail'' --- it
is ``if this fails, the failure voids my parent too.''
\end{definition}

The practical reading: make fields non-null when a client cannot
meaningfully render anything without them, and nullable when partial
degradation is useful. A product page can render without a margin; a
payment screen cannot render without an amount. Our reference schema
encodes exactly this: \texttt{marginPercent} is \texttt{Int!}, so its
failure (cost data missing for one discontinued product) nulls the
enclosing node while sibling products survive ---
Listing~\ref{lst:ch12-tests} proves both halves of that sentence.

\subsection{Schema Definition Language and Introspection}
\label{sec:ch12-sdl}

The schema definition language (SDL) is the human-readable notation for
the type system. The catalog API's core, as SDL, is:

\begin{minted}{graphql}
type Query {
  products(first: Int = 10, after: String, status: ProductStatus): ProductConnection
  product(productId: Int!): Product
}

type Product {
  id: Int!
  sku: String!
  name: String!
  unitPriceCents: Int!
  stock: Int!
  discontinued: Boolean!
  category: Category
  reviews: [Review!]!
  marginPercent: Int!
}

type Mutation {
  addReview(input: AddReviewInput!): AddReviewResult!
}

union AddReviewResult = AddReviewSuccess | AddReviewError
\end{minted}

SDL is one authoring style; Strawberry's is \emph{code-first}, where Python
classes generate the SDL. The choice matters less than the invariant:
\textbf{the schema is the single source of truth}, and anything claiming
to describe the API must be derived from it.

Introspection is the schema describing \emph{itself} at runtime: the
meta-fields \texttt{\_\_schema} and \texttt{\_\_type} expose every type,
field, and argument, which is what powers GraphiQL's autocompletion and
every code generator. It is also a complete map of your attack surface,
served fresh to anyone who asks.

\begin{warningbox}
Introspection is not an authentication bypass --- it reveals only types,
not data --- but it hands an adversary your entire field inventory,
including the internal field you forgot carried something sensitive.
Production policy: introspection \textbf{off} for unauthenticated traffic
in production, \textbf{on} in development and staging, and gate it in the
validation layer (an AST check for \texttt{\_\_schema}/\texttt{\_\_type}),
never by substring matching --- Listing~\ref{lst:ch12-cost} implements the
AST gate that aliases and fragments cannot evade.
\end{warningbox}

\subsection{Operations: Queries, Mutations, Subscriptions}
\label{sec:ch12-operations}

A \textbf{query} is a read with no observable side effects; a
\textbf{mutation} is a write whose top-level fields execute \emph{in
series}, in order --- the only ordering guarantee the specification makes,
and the reason multi-step writes belong in one mutation's sequential
fields rather than in parallel mutations; a \textbf{subscription} is a
query plus a stream: the server pushes a response each time an event
invalidates the selection, typically over WebSockets (the
\texttt{graphql-ws} protocol) or server-sent events.

Operations accept \textbf{variables}: typed, separately-serialized
arguments declared in the operation header (\texttt{query Page(\$first:
Int!)}). Variables are the only correct way to pass client input --- a
query is a parse tree, and variables travel in their own JSON object, so
\emph{variable injection is structurally impossible}: there is no string
to escape from. \textbf{Aliases} (\texttt{active: products(...)}) rename
result keys, permitting the same field twice with different arguments.
\textbf{Fragments} name reusable selection sets; the directives
\texttt{@include(if:)} and \texttt{@skip(if:)} condition fields on
variables, letting one cached operation serve several UI states.

\begin{minted}{graphql}
query CatalogPage($first: Int!, $after: String, $withReviews: Boolean!) {
  products(first: $first, after: $after) {
    edges {
      cursor
      node {
        ...ProductCard
        reviews @include(if: $withReviews) { rating }
      }
    }
    pageInfo { hasNextPage endCursor }
  }
}

fragment ProductCard on Product {
  sku
  name
  unitPriceCents
}
\end{minted}

\subsection{The Execution Pipeline and the Wire Shape}
\label{sec:ch12-pipeline}

Every GraphQL request walks the same three phases, and each phase has its
own failure channel (Figure~\ref{fig:ch12-pipeline}). \textbf{Parse} turns
the document text into an AST; failure is a syntax error, HTTP 400, before
any schema involvement. \textbf{Validate} checks the AST against the schema
--- every field exists, every argument type-checks, fragments terminate ---
failure again pre-execution. This is also where cost control belongs:
depth and complexity are properties of the \emph{document}, so an
over-budget operation must be rejected before a resolver runs, or the
meter is already spinning. \textbf{Execute} resolves fields; failures here
are field errors, folded into the response under the null-propagation rule.

\begin{figure}[h]
  \centering
  \begin{tikzpicture}[
      node distance=0.55cm,
      phase/.style={rectangle, draw=packtgray, thick, fill=blue!5,
        minimum width=2.5cm, minimum height=1.0cm, align=center},
      gate/.style={rectangle, draw=packtorange, thick, fill=packtorange!10,
        minimum width=3.2cm, minimum height=0.9cm, align=center},
      io/.style={rectangle, draw=packtgray!70, fill=lightgray,
        minimum width=3.4cm, minimum height=0.9cm, align=center},
      flow/.style={->, thick, packtgray},
      reject/.style={->, thick, packtorange, dashed},
    ]
    \node[io] (req) {request envelope\\\texttt{query}, \texttt{variables}, \texttt{operationName}};
    \node[phase, below=of req] (parse) {\textbf{Parse}\\text $\to$ AST};
    \node[gate, below=of parse] (budget) {\textbf{Cost gate}\\depth, complexity,\\persisted allow-list};
    \node[phase, below=of budget] (validate) {\textbf{Validate}\\AST vs.\ schema};
    \node[phase, below=of validate] (execute) {\textbf{Execute}\\resolvers, DataLoaders};
    \node[io, below=of execute] (resp) {response envelope\\\texttt{data} + \texttt{errors}};
    \node[io, right=1.3cm of parse] (r1) {HTTP 400\\\texttt{GRAPHQL\_PARSE\_FAILED}};
    \node[io, right=1.3cm of budget] (r2) {HTTP 400\\\texttt{DEPTH\_LIMIT\_EXCEEDED}};
    \node[io, right=1.3cm of validate] (r3) {HTTP 400\\validation errors};
    \node[io, right=1.3cm of execute, align=center] (r4) {HTTP 200, partial\\\texttt{errors} with \texttt{path}};
    \draw[flow] (req) -- (parse);
    \draw[flow] (parse) -- (budget);
    \draw[flow] (budget) -- (validate);
    \draw[flow] (validate) -- (execute);
    \draw[flow] (execute) -- (resp);
    \draw[reject] (parse) -- (r1);
    \draw[reject] (budget) -- (r2);
    \draw[reject] (validate) -- (r3);
    \draw[reject] (execute) -- (r4);
    \node[note, below=0.25cm of resp] {phases left of the dashed rejects never touch a resolver};
  \end{tikzpicture}
  \caption{Request lifecycle. Static cost control sits between parse and
           validate: cheap to compute, before any field executes.}
  \label{fig:ch12-pipeline}
\end{figure}

The wire shape is austere. The request is a JSON object with
\texttt{query}, optional \texttt{variables}, optional
\texttt{operationName}, and optional \texttt{extensions}; over HTTP it is
almost always POST \cite{rfc9110}. The response is a JSON object with
\texttt{data} (possibly partial, possibly \texttt{null}) and, on field
errors, an \texttt{errors} array whose entries carry \texttt{message},
\texttt{path}, \texttt{locations}, and the open extension point
\texttt{extensions} --- where \texttt{code} belongs, because clients need
machine-readable failure kinds, not prose.

\begin{minted}{json}
{
  "data": { "product": null },
  "errors": [
    {
      "message": "cost data unavailable for product",
      "path": ["product", "marginPercent"],
      "extensions": { "code": "COST_DATA_UNAVAILABLE" }
    }
  ]
}
\end{minted}

Read that response carefully, because it breaks two REST instincts at
once: the HTTP status is 200 \emph{and} the operation failed, and the
failure voided exactly one field's parent, not the request. A client that
checks only \texttt{resp.status\_code} will render garbage; a client that
treats any \texttt{errors} entry as total failure wastes the partial data
the server worked to produce. Both behaviors are bugs.

\subsection{Resolver Architecture and the N+1 Problem}
\label{sec:ch12-resolvers}

Execution is field-at-a-time: for each field in the selection set, the
executor calls that field's \emph{resolver} with the parent value, the
arguments, and the \emph{context} --- a per-request bag holding
repositories, the authenticated principal, and (critically) DataLoaders.
This design is what makes GraphQL composable: the \texttt{category}
resolver does not know or care whether the parent product came from
\texttt{products} or \texttt{product}. It is also what makes GraphQL
dangerous, because resolution composes \emph{multiplicatively} with list
sizes.

\begin{examplebox}
The query \texttt{products(first: 5) \{ edges \{ node \{ name category
\{ name \} reviews \{ rating \} \} \} \} } looks like one request. With
naive resolvers --- \texttt{category} calls
\texttt{repo.category\_by\_id}, \texttt{reviews} calls
\texttt{repo.reviews\_for\_product} --- it is $1 + 5 + 5 = 11$ store
reads: one for the page, then one \emph{per product} per relationship
field. At \texttt{first: 50} it is 101 reads. The $N{+}1$ is not a bug in
a resolver; it is the default physics of field-level resolution.
Listing~\ref{lst:ch12-nplus1} measures exactly these counts.
\end{examplebox}

DataLoader collapses the fan-out by exploiting the executor's scheduling:
all product nodes' \texttt{category} resolvers run within the same event
loop tick, so each one, instead of querying, \emph{enqueues its key} and
awaits. At the tick boundary the loader calls the batch function once with
the deduplicated keys --- \texttt{[1, 1, 2, 2, 3]} becomes
\texttt{categories\_by\_ids([1, 2, 3])} --- and distributes results back
to the waiting futures by position. Eleven reads become three: the page,
one batched category read, one batched review read
(Figure~\ref{fig:ch12-dataloader}).

\begin{figure}[h]
  \centering
  \begin{tikzpicture}[
      lifeline/.style={rectangle, draw=packtgray, thick, fill=blue!5,
        minimum width=2.6cm, minimum height=0.8cm, align=center},
      msg/.style={->, thick, packtgray},
      ret/.style={->, thick, dashed, packtgray!80},
      lbl/.style={font=\footnotesize, align=center, midway, above},
    ]
    \node[lifeline] (r1) {resolver: product 1};
    \node[lifeline, right=0.7cm of r1] (r2) {resolver: product 2};
    \node[lifeline, right=0.7cm of r2] (dl) {CategoryLoader};
    \node[lifeline, right=0.7cm of dl] (db) {repository};
    \draw[packtgray!40, dashed] (r1.south) -- ++(0,-4.6);
    \draw[packtgray!40, dashed] (r2.south) -- ++(0,-4.6);
    \draw[packtgray!40, dashed] (dl.south) -- ++(0,-4.6);
    \draw[packtgray!40, dashed] (db.south) -- ++(0,-4.6);
    \draw[msg] ([yshift=-0.5cm]r1.south) -- ([yshift=-0.5cm]dl.south)
      node[lbl] {\texttt{load(1)} --- enqueue, await};
    \draw[msg] ([yshift=-1.1cm]r2.south) -- ([yshift=-1.1cm]dl.south)
      node[lbl] {\texttt{load(1)} --- cache hit, same future};
    \draw[msg] ([yshift=-1.9cm]dl.south) -- ([yshift=-1.9cm]db.south)
      node[lbl] {tick boundary: \texttt{categories\_by\_ids([1])} --- 1 query};
    \draw[ret] ([yshift=-2.9cm]db.south) -- ([yshift=-2.9cm]dl.south)
      node[lbl] {\texttt{[Category(1)]}};
    \draw[ret] ([yshift=-3.7cm]dl.south) -- ([yshift=-3.7cm]r1.south)
      node[lbl] {resolve futures};
    \draw[ret] ([yshift=-4.4cm]dl.south) -- ([yshift=-4.4cm]r2.south);
  \end{tikzpicture}
  \caption{DataLoader batching: per-field loads within one tick collapse
           into a single keyed batch read; the per-request cache dedupes
           repeated keys.}
  \label{fig:ch12-dataloader}
\end{figure}

Two rules keep DataLoader correct. \textbf{One loader instance per
request}: its cache is request-scoped identity memoization; sharing a
loader across requests serves stale data across user boundaries, which is
not a performance bug but a security incident. \textbf{The batch function
is the only place queries happen}: it must return results positionally
aligned with the keys, with \texttt{None} for misses --- never raise for
one bad key, or every awaiting field in the tick fails together.

\subsection{Pagination: The Relay Connection Spec}
\label{sec:ch12-pagination}

Chapter~8 established that value-based (cursor) pagination beats
position-based (offset) pagination; the Relay connection spec is the
GraphQL community's standard encoding of that idea. A paginated field
returns a \emph{connection} object: \texttt{edges} (each with an opaque
\texttt{cursor} and a \texttt{node}), \texttt{pageInfo} (with
\texttt{hasNextPage}, \texttt{hasPreviousPage}, \texttt{startCursor},
\texttt{endCursor}), and forward arguments \texttt{first}/\texttt{after}
(backward: \texttt{last}/\texttt{before}). The conventions that matter in
production: cursors are \emph{opaque to clients} (ours are
Base64\texttt{("product:<id>")}, decodable server-side, forged ones
rejected with \texttt{BAD\_CURSOR}); \texttt{hasNextPage} is computed with
the fetch-$p{+}1$ trick, never with a \texttt{COUNT}; and every list a
client can size needs a server-side maximum --- our schema caps
\texttt{first} at 50 and rejects out-of-range values at the resolver.

\subsection{The Error Model}
\label{sec:ch12-errors}

GraphQL has exactly two error channels, and choosing between them is
design. \textbf{Top-level errors} (the \texttt{errors} array) are for
\emph{exceptional} failure: the resolver could not produce the field ---
downstream timeout, missing cost data, invalid cursor. They compose with
null propagation and partial data. \textbf{Typed mutation errors} model
\emph{expected} failure --- validation, business-rule rejection --- as
data: the mutation returns a union of success and error payloads, and the
client selects with inline fragments:

\begin{minted}{graphql}
mutation Add($input: AddReviewInput!) {
  addReview(input: $input) {
    __typename
    ... on AddReviewSuccess { review { id rating } }
    ... on AddReviewError { code message }
  }
}
\end{minted}

The distinction is operational, not stylistic. Top-level errors page the
on-call engineer; typed errors inform the end user. A rating of 6 is not
an incident, and a database connection failure is not a form-validation
message. Chapter~7's discipline --- machine-readable codes, actionable
messages --- applies to both channels via \texttt{extensions.code} on the
one hand and the \texttt{code} field on the other.

\subsection{API Cost Control}
\label{sec:ch12-cost-control}

Because the client chooses the work, an unguarded GraphQL endpoint has an
unbounded worst case. The controls, in the order they should fire:

\textbf{Depth limiting} rejects documents whose longest field path exceeds
a budget (ours: 7). Depth is the crudest measure but catches the classic
abuse --- cyclic traversals like \texttt{friends \{ friends \{ friends
\ldots\ \} \}} --- before they cost anything.

\textbf{Complexity analysis} weights fields so that list-returning fields
multiply their subtree's cost by the requested page size:
$\mathrm{cost}(f) = 1 + m(f) \cdot \sum_{c \in \mathrm{children}(f)}
\mathrm{cost}(c)$, where $m(f)$ is the field's \texttt{first}/\texttt{last}
argument, a conservative default when absent, and 1 for scalar fields. The
estimator runs on the parsed AST (Listing~\ref{lst:ch12-cost}), so a
\texttt{first: 50} god-query is rejected with
\texttt{COMPLEXITY\_LIMIT\_EXCEEDED} while a \texttt{first: 2} variant of
the same shape sails through --- the suite proves both.

\textbf{Timeouts} bound what static analysis cannot predict; complexity
estimates are adversarial games, and a resolver deadline is the backstop
for the estimation errors you did not know you had.

\textbf{Operation allow-lists} invert the trust model entirely: in
production the server executes only pre-registered documents, identified
by hash. Nothing else parses, so depth and complexity limits become
defense in depth rather than the perimeter. \textbf{Automatic Persisted
Queries} (APQ) is the ergonomic half of that idea: the client sends only
\texttt{extensions.persistedQuery.sha256Hash}; on
\texttt{PersistedQueryNotFound} (HTTP 200 with an error entry, per the
protocol) it retries once with the full query plus hash; the server
verifies the hash matches the text --- \texttt{PersistedQueryMismatch}
otherwise --- registers the document, and executes. Subsequent requests
carry 64 bytes instead of kilobytes, and a CDN can cache hash-keyed GETs.
Listing~\ref{lst:ch12-client} implements both halves of the handshake.

\subsection{Security: The GraphQL-Specific Attack Surface}
\label{sec:ch12-security}

GraphQL inherits every web-API threat and adds a few of its own.

\textbf{Introspection abuse} --- covered in
Section~\ref{sec:ch12-sdl}: gate it at the validation layer.
\textbf{Query-based DoS} --- deep, wide, or aliased-multiplied documents;
mitigated by the cost gate above. Note the amplifier most teams miss:
\textbf{batching abuse}. GraphQL-over-HTTP conventionally accepts an
\emph{array} of operations per request, so one HTTP request can carry a
hundred god-queries, and request-count rate limiting sees ``1 request.''
Mitigations: disable array batching if your clients do not need it, cap
the array length, and --- the durable fix --- \textbf{rate-limit by
estimated cost, not by request count}, summing complexity across the
batch. \textbf{Injection}: variable injection into the \emph{query
document} is structurally impossible (variables are a separate JSON
object), but \emph{resolver-side} injection is exactly as real as in any
stack --- a resolver that string-concatenates a variable into SQL is a SQL
injection wearing a typed costume. Parameterized queries remain
non-negotiable \cite{kleppmann2017ddia}. \textbf{Field-level
authorization} moves the perimeter into resolvers: the schema advertises
\texttt{salaryHistory} to everyone, so authorization must live at the
field, via context-carrying checks --- and union payloads let you return
\texttt{NOT\_AUTHORIZED} as data where partial visibility is legitimate.
Finally, \textbf{disable the GraphiQL IDE in production}: it is a
developer tool that leaks the endpoint's shape and invites interactive
probing; its presence is a config flag, and that flag belongs in the same
policy file as the introspection gate.

\subsection{GraphQL vs.\ REST vs.\ gRPC: An Honest Matrix}
\label{sec:ch12-comparison}

\begin{table}[h]
  \centering
  \small
  \begin{tabularx}{\textwidth}{l X X X}
    \toprule
    \textbf{Dimension} & \textbf{REST (HTTP+JSON)} & \textbf{GraphQL} & \textbf{gRPC} \\
    \midrule
    Contract & OpenAPI, retrofitted & Schema, intrinsic \cite{graphqlspec} & Protobuf, intrinsic \\
    Response shape & Server-fixed & Client-composed & Server-fixed \\
    HTTP caching & Native: GET + CDN + ETag & POST defeats it; APQ+GET partially restores & No (HTTP/2 frames) \\
    Over/under-fetching & Common & Eliminated & Eliminated \\
    File uploads & Native multipart & Multipart spec, awkward & Streaming, good \\
    Versioning & \texttt{/v2} or headers & Continuous evolution: deprecate, never break & Field numbers, additive \\
    Worst-case cost & Bounded per endpoint & Unbounded without cost gate & Bounded per method \\
    Tooling gravity & Ubiquitous & Strong for frontends & Strong for service mesh \\
    Fit & Public APIs, webhooks, CDNs & Product graphs, many clients, one backend & Internal low-latency RPC \\
    \bottomrule
  \end{tabularx}
  \caption{Decision matrix. No column is ``best''; each optimizes a
           different constraint.}
  \label{tab:ch12-matrix}
\end{table}

Three rows deserve emphasis. \textbf{Cacheability}: REST's use of GET with
status codes and validators is the largest operational advantage in the
table --- a CDN can absorb a product launch; a GraphQL POST endpoint
cannot, and the mitigations (APQ over GET, response-caching extensions)
are exactly that: mitigations. \textbf{Versioning philosophy}: GraphQL's
orthodoxy is \emph{continuous evolution} --- add fields freely, deprecate
with \texttt{@deprecated}, never remove --- which trades REST's explicit
\texttt{/v2} flag days for an eternal compatibility obligation backed by
usage telemetry on every field. Neither is cheaper; they bill different
teams. \textbf{The request topology} differs structurally
(Figure~\ref{fig:ch12-topology}): REST spreads a screen's data need across
$k$ endpoints and $k$ round trips; GraphQL concentrates it in one request
whose cost you must now govern.

\begin{figure}[h]
  \centering
  \begin{tikzpicture}[
      box/.style={rectangle, draw=packtgray, thick, fill=blue!5,
        minimum width=1.9cm, minimum height=0.8cm, align=center},
      ep/.style={rectangle, draw=packtgray, fill=lightgray,
        minimum width=1.9cm, minimum height=0.7cm, align=center,
        font=\footnotesize},
      flow/.style={->, thick, packtgray},
      ttl/.style={font=\bfseries\small},
    ]
    % REST side
    \node[ttl] at (2.2, 0.4) {REST: $k$ endpoints, $k$ round trips};
    \node[box] (c1) at (0.2, -0.6) {mobile client};
    \node[ep] (e1) at (4.3, 0.0) {\texttt{GET /products}};
    \node[ep] (e2) at (4.3, -0.9) {\texttt{GET /products/1/category}};
    \node[ep] (e3) at (4.3, -1.8) {\texttt{GET /products/1/reviews}};
    \draw[flow] (c1) -- (e1);
    \draw[flow] (c1) -- (e2);
    \draw[flow] (c1) -- (e3);
    % GraphQL side
    \node[ttl] at (9.6, 0.4) {GraphQL: 1 endpoint, 1 request};
    \node[box] (c2) at (7.6, -0.6) {mobile client};
    \node[ep, minimum height=1.6cm] (g1) at (11.6, -0.6)
      {\texttt{POST /graphql}\\one composed operation};
    \draw[flow] (c2) -- (g1);
    \node[note] at (9.6, -2.0) {cost governance moves server-side};
  \end{tikzpicture}
  \caption{Request topology: REST distributes a screen's data need across
           endpoints; GraphQL concentrates it into one operation whose
           cost the server must bound.}
  \label{fig:ch12-topology}
\end{figure}

\subsection{Federation Primer (Scope Note)}
\label{sec:ch12-federation}

When one schema is owned by fifteen teams, the monolith reappears wearing
a schema's clothes. Federation splits the graph into \emph{subgraphs} ---
each a standalone GraphQL service owning its types --- that a
\emph{gateway} composes into one supergraph; \emph{entities} are types
with a key (say, \texttt{Product @key(fields: "id")}) that several
subgraphs can extend, so the reviews service adds
\texttt{Product.reviews} without owning products. The gateway plans each
operation across subgraphs and stitches the results. That is the entire
scope of this chapter's treatment: federation is an organizational
scaling tool with its own failure modes (composition conflicts, gateway as
a chokepoint, entity-resolution $N{+}1$s \emph{across services}), and it
deserves the same rigor this chapter gave the single graph before you
adopt it.

%% ============================================================================
\section{Production-Ready Code Implementation}
\label{sec:ch12-code}

The reference implementation is a six-file lab, deliberately offline and
deterministic. \textbf{Dependencies.} We use
\textbf{Strawberry} rather than Ariadne or Graphene: it is code-first
(Python type hints \emph{are} the schema, so the listing teaches the type
system instead of an SDL mapping layer), it is actively maintained on
modern \texttt{graphql-core}, its DataLoader and FastAPI integration are
first-class, and --- decisive for a textbook --- its resolvers are plain
async Python that reads like the specification. Ariadne (SDL-first) and
Graphene (the incumbent) are surveyed in the Dependencies Appendix.

\textbf{Layout.} \texttt{catalog\_repository.py} (seeded store with query
counter), \texttt{catalog\_schema.py} (schema, connections, mutations,
DataLoaders), \texttt{cost\_control.py} (depth/complexity/APQ/HTTP entry
point), \texttt{server.py} (FastAPI host), \texttt{graphql\_client.py}
(httpx client with APQ), \texttt{test\_catalog.py} (21 in-process tests).
Windows paths below assume the lab lives at
\texttt{C:\textbackslash{}training\textbackslash{}graphql\_lab}.

\subsection{Listing 1: Instrumented Repository}
\label{subsec:ch12-code-repository}

The repository (Listing~\ref{lst:ch12-repository}) is the lab's ground
truth: five synthetic products across three categories with five reviews,
one of them (product~4, the discontinued sensor array) deliberately missing
cost data to drive the null-propagation demonstration. Every read path
increments \texttt{query\_count}, which converts the $N{+}1$ problem from
anecdote into an assertion a test can make.

\begin{minted}[caption={Instrumented in-memory repository (\texttt{catalog\_repository.py}).},label={lst:ch12-repository}]{python}
"""In-memory repository for the Northwind Robotics catalog.

Synthetic, deterministic, offline: the same seed data on every run.
The repository exposes an instrumentation counter so the N+1 problem is
measurable rather than anecdotal.
"""

from __future__ import annotations

from dataclasses import dataclass


@dataclass(frozen=True)
class CategoryRow:
    id: int
    name: str
    description: str


@dataclass(frozen=True)
class ProductRow:
    id: int
    sku: str
    name: str
    category_id: int
    unit_price_cents: int
    unit_cost_cents: int | None  # None = cost data missing (drives null-propagation demo)
    stock: int
    discontinued: bool


@dataclass(frozen=True)
class ReviewRow:
    id: int
    product_id: int
    rating: int
    headline: str


_CATEGORIES: tuple[CategoryRow, ...] = (
    CategoryRow(1, "Actuators", "Linear and rotary motion components"),
    CategoryRow(2, "Sensors", "LIDAR, vision, and tactile sensing"),
    CategoryRow(3, "Controllers", "Real-time motor and fleet controllers"),
)

_PRODUCTS: tuple[ProductRow, ...] = (
    ProductRow(1, "NWR-ACT-100", "Servo Linear Actuator 100mm", 1, 459_00, 210_00, 34, False),
    ProductRow(2, "NWR-ACT-200", "Rotary Torque Actuator 12Nm", 1, 899_00, 401_00, 12, False),
    ProductRow(3, "NWR-SEN-310", "Solid-State LIDAR Unit", 2, 1_299_00, 620_00, 8, False),
    ProductRow(4, "NWR-SEN-410", "Tactile Gripper Sensor Array", 2, 349_00, None, 0, True),
    ProductRow(5, "NWR-CTL-900", "Fleet Motion Controller X9", 3, 2_499_00, 1_150_00, 5, False),
)

_REVIEWS: tuple[ReviewRow, ...] = (
    ReviewRow(1, 1, 5, "Silky smooth travel"),
    ReviewRow(2, 1, 4, "Solid, runs warm under load"),
    ReviewRow(3, 2, 5, "Best torque-per-dollar we tested"),
    ReviewRow(4, 3, 4, "Great range, tricky mounting"),
    ReviewRow(5, 5, 5, "Runs our whole warehouse fleet"),
)


class CatalogRepository:
    """In-memory catalog store with a query counter for N+1 measurement."""

    def __init__(self) -> None:
        self._categories = {c.id: c for c in _CATEGORIES}
        self._products = {p.id: p for p in _PRODUCTS}
        self._reviews: dict[int, list[ReviewRow]] = {}
        for review in _REVIEWS:
            self._reviews.setdefault(review.product_id, []).append(review)
        self._next_review_id = max(r.id for r in _REVIEWS) + 1
        self.query_count = 0  # incremented on every store read

    def reset_counter(self) -> None:
        self.query_count = 0

    # --- single-record reads (the naive path uses these in a loop) ---

    def product_by_id(self, product_id: int) -> ProductRow | None:
        self.query_count += 1
        return self._products.get(product_id)

    def category_by_id(self, category_id: int) -> CategoryRow | None:
        self.query_count += 1
        return self._categories.get(category_id)

    def reviews_for_product(self, product_id: int) -> list[ReviewRow]:
        self.query_count += 1
        return list(self._reviews.get(product_id, []))

    # --- batched reads (the DataLoader path calls these once per tick) ---

    def categories_by_ids(self, ids: list[int]) -> list[CategoryRow | None]:
        self.query_count += 1
        return [self._categories.get(i) for i in ids]

    def reviews_for_products(self, ids: list[int]) -> list[list[ReviewRow]]:
        self.query_count += 1
        return [list(self._reviews.get(i, [])) for i in ids]

    # --- connection page read ---

    def products_page(
        self, after_id: int, first: int
    ) -> tuple[list[ProductRow], bool]:
        """Return up to ``first`` products with id > after_id, plus has-next."""
        self.query_count += 1
        ordered = sorted(self._products.values(), key=lambda p: p.id)
        page = [p for p in ordered if p.id > after_id][: first + 1]
        return page[:first], len(page) > first

    def total_products(self) -> int:
        self.query_count += 1
        return len(self._products)

    # --- mutation ---

    def add_review(self, product_id: int, rating: int, headline: str) -> ReviewRow | None:
        """Insert a review; returns None when the product does not exist."""
        if product_id not in self._products:
            return None
        self._next_review_id += 1
        row = ReviewRow(self._next_review_id, product_id, rating, headline)
        self._reviews.setdefault(product_id, []).append(row)
        return row
\end{minted}

Note the deliberate asymmetry: single-record reads
(\texttt{category\_by\_id}, \texttt{reviews\_for\_product}) exist alongside
batched reads (\texttt{categories\_by\_ids},
\texttt{reviews\_for\_products}). The naive resolvers use the first set;
the DataLoaders use the second. Same store, same data --- the only
difference is \emph{how many round trips the resolvers ask for}.

\subsection{Listing 2: The Strawberry Schema}
\label{subsec:ch12-code-schema}

Listing~\ref{lst:ch12-schema} is the chapter in one file. Code-first type
definitions with Python hints generating the SDL of
Section~\ref{sec:ch12-sdl}; a \texttt{Context} constructed \emph{once per
request} carrying the repository and two DataLoaders; a Relay-style
connection with opaque Base64 cursors; a mutation returning a union of
\texttt{AddReviewSuccess} and \texttt{AddReviewError} so expected failures
are data; and \texttt{margin\_percent} --- a non-null field that raises a
\texttt{GraphQLError} carrying \texttt{extensions.code} when cost data is
missing, so the tests can watch null propagation happen.

\begin{minted}[caption={Catalog schema with connections, union-typed mutations, and DataLoaders (\texttt{catalog\_schema.py}).},label={lst:ch12-schema}]{python}
"""Strawberry schema for the Northwind Robotics catalog.

Demonstrates: object types, enums, input types, union-typed mutation
results, Relay-style connections, context injection, per-request
DataLoaders, field-level errors with extensions, and null propagation.
"""

from __future__ import annotations

import base64
import enum
from typing import Annotated, Optional

import strawberry
from graphql import GraphQLError
from strawberry.dataloader import DataLoader
from strawberry.types import Info

from catalog_repository import CatalogRepository, CategoryRow, ProductRow, ReviewRow

MAX_PAGE_SIZE = 50


# ---------------------------------------------------------------------------
# Enum and context
# ---------------------------------------------------------------------------


@strawberry.enum
class ProductStatus(enum.Enum):
    ACTIVE = "active"
    DISCONTINUED = "discontinued"


class Context:
    """Per-request context: one repository, one set of DataLoaders.

    NOT a Strawberry type: it is constructed once per request and handed
    to resolvers via ``info.context``. ``use_loaders`` exists only so the
    N+1 demo can switch resolvers between naive and batched reads.
    """

    def __init__(self, repo: CatalogRepository, use_loaders: bool = True) -> None:
        self.repo = repo
        self.use_loaders = use_loaders
        self.category_loader = make_category_loader(repo)
        self.review_loader = make_review_loader(repo)


def get_context() -> Context:
    return Context(CatalogRepository())


# ---------------------------------------------------------------------------
# Object types
# ---------------------------------------------------------------------------


@strawberry.type
class Category:
    id: int
    name: str
    description: str


@strawberry.type
class Review:
    id: int
    rating: int
    headline: str


@strawberry.type
class Product:
    id: int
    sku: str
    name: str
    unit_price_cents: int
    stock: int
    discontinued: bool
    _category_id: strawberry.Private[int]

    @strawberry.field
    async def category(self, info: Info[Context, None]) -> Optional[Category]:
        ctx = info.context
        if ctx.use_loaders:
            row = await ctx.category_loader.load(self._category_id)
        else:  # naive path: one store read per product (the N in N+1)
            row = ctx.repo.category_by_id(self._category_id)
        if row is None:
            return None
        return Category(id=row.id, name=row.name, description=row.description)

    @strawberry.field
    async def reviews(self, info: Info[Context, None]) -> list[Review]:
        ctx = info.context
        if ctx.use_loaders:
            rows = await ctx.review_loader.load(self.id)
        else:  # naive path
            rows = ctx.repo.reviews_for_product(self.id)
        return [Review(id=r.id, rating=r.rating, headline=r.headline) for r in rows]

    @strawberry.field
    def margin_percent(self, info: Info[Context, None]) -> int:
        """NON-NULL field that errors when cost data is missing.

        Raising here nulls the entire Product (null propagation), because
        the schema promises this field is never null.
        """
        row = info.context.repo.product_by_id(self.id)
        if row is None or row.unit_cost_cents is None:
            raise GraphQLError(
                "cost data unavailable for product",
                extensions={"code": "COST_DATA_UNAVAILABLE"},
            )
        margin = row.unit_price_cents - row.unit_cost_cents
        return round(100 * margin / row.unit_price_cents)


def _to_product(row: ProductRow) -> Product:
    return Product(
        id=row.id,
        sku=row.sku,
        name=row.name,
        unit_price_cents=row.unit_price_cents,
        stock=row.stock,
        discontinued=row.discontinued,
        _category_id=row.category_id,
    )


# ---------------------------------------------------------------------------
# Relay-style connection
# ---------------------------------------------------------------------------


def encode_cursor(product_id: int) -> str:
    raw = f"product:{product_id}".encode("ascii")
    return base64.urlsafe_b64encode(raw).decode("ascii")


def decode_cursor(cursor: str) -> int:
    try:
        raw = base64.urlsafe_b64decode(cursor.encode("ascii")).decode("ascii")
        prefix, sep, digits = raw.partition(":")
        if prefix != "product" or sep == "":
            raise ValueError("bad prefix")
        return int(digits)
    except (ValueError, UnicodeDecodeError) as exc:
        raise GraphQLError(
            "invalid cursor", extensions={"code": "BAD_CURSOR"}
        ) from exc


@strawberry.type
class PageInfo:
    has_next_page: bool
    has_previous_page: bool
    start_cursor: Optional[str]
    end_cursor: Optional[str]


@strawberry.type
class ProductEdge:
    cursor: str
    node: Optional[Product]  # nullable: a nulled node must not void the page


@strawberry.type
class ProductConnection:
    edges: list[ProductEdge]
    page_info: PageInfo

    @strawberry.field
    def total_count(self, info: Info[Context, None]) -> int:
        # Lazy: only costs a store read when the client asks for it.
        return info.context.repo.total_products()


# ---------------------------------------------------------------------------
# Mutation payloads: typed errors via unions
# ---------------------------------------------------------------------------


@strawberry.input
class AddReviewInput:
    product_id: int
    rating: int
    headline: str


@strawberry.type
class AddReviewSuccess:
    review: Review


@strawberry.type
class AddReviewError:
    message: str
    code: str


AddReviewResult = Annotated[
    AddReviewSuccess | AddReviewError, strawberry.union("AddReviewResult")
]


# ---------------------------------------------------------------------------
# Root operations
# ---------------------------------------------------------------------------


@strawberry.type
class Query:
    @strawberry.field
    async def products(
        self,
        info: Info[Context, None],
        first: int = 10,
        after: Optional[str] = None,
        status: Optional[ProductStatus] = None,
    ) -> Optional[ProductConnection]:  # nullable root: errors stay partial
        if not 1 <= first <= MAX_PAGE_SIZE:
            raise GraphQLError(
                f"first must be between 1 and {MAX_PAGE_SIZE}",
                extensions={"code": "BAD_PAGINATION_ARGUMENT"},
            )
        after_id = decode_cursor(after) if after else 0
        repo = info.context.repo
        rows, has_next = repo.products_page(after_id, first)
        if status is ProductStatus.ACTIVE:
            rows = [r for r in rows if not r.discontinued]
        elif status is ProductStatus.DISCONTINUED:
            rows = [r for r in rows if r.discontinued]
        edges = [
            ProductEdge(cursor=encode_cursor(r.id), node=_to_product(r)) for r in rows
        ]
        return ProductConnection(
            edges=edges,
            page_info=PageInfo(
                has_next_page=has_next,
                has_previous_page=after_id > 0,
                start_cursor=edges[0].cursor if edges else None,
                end_cursor=edges[-1].cursor if edges else None,
            ),
        )

    @strawberry.field
    async def product(self, info: Info[Context, None], product_id: int) -> Optional[Product]:
        row = info.context.repo.product_by_id(product_id)
        return _to_product(row) if row else None


@strawberry.type
class Mutation:
    @strawberry.mutation
    def add_review(
        self, info: Info[Context, None], input: AddReviewInput
    ) -> AddReviewResult:
        if not 1 <= input.rating <= 5:
            return AddReviewError(
                message="rating must be between 1 and 5",
                code="RATING_OUT_OF_RANGE",
            )
        headline = input.headline.strip()
        if not headline:
            return AddReviewError(
                message="headline must not be empty", code="EMPTY_HEADLINE"
            )
        row = info.context.repo.add_review(input.product_id, input.rating, headline)
        if row is None:
            return AddReviewError(
                message=f"product {input.product_id} not found",
                code="PRODUCT_NOT_FOUND",
            )
        return AddReviewSuccess(
            review=Review(id=row.id, rating=row.rating, headline=row.headline)
        )


# ---------------------------------------------------------------------------
# DataLoaders: one instance per request, batch + per-request cache
# ---------------------------------------------------------------------------


def make_category_loader(repo: CatalogRepository) -> DataLoader:
    async def batch(keys: list[int]) -> list[Optional[CategoryRow]]:
        # ONE store read per tick, for all keys queued this tick.
        return repo.categories_by_ids(keys)

    return DataLoader(load_fn=batch)


def make_review_loader(repo: CatalogRepository) -> DataLoader:
    async def batch(keys: list[int]) -> list[list[ReviewRow]]:
        return repo.reviews_for_products(keys)

    return DataLoader(load_fn=batch)


schema = strawberry.Schema(query=Query, mutation=Mutation)
\end{minted}

Two details deserve a second look. \texttt{total\_count} is a resolver, not
a dataclass field: it costs a store read only when the client actually
selects \texttt{totalCount}, which is the whole philosophy of GraphQL
cost accounting in miniature --- you pay for what you select, and the
schema author decides what each selection costs. And
\texttt{use\_loaders} exists \emph{only} so the next listing can flip the
same resolvers between naive and batched store access without changing
the query; it is a lab instrument, not a production pattern.

\subsection{Listing 3: The N+1 Demonstration}
\label{subsec:ch12-code-nplus1}

Listing~\ref{lst:ch12-nplus1} executes one operation --- five products
with category and reviews --- twice, counting store reads. The output is
the chapter's most important four lines:

\begin{minted}{text}
GraphQL operation asks for 5 products + category + reviews
  naive resolvers   :  11 store queries (1 page + 5 category + 5 review reads)
  DataLoader path   :   3 store queries (1 page + 1 batched category read + 1 batched review read)
N+1 reproduced and collapsed. OK
\end{minted}

\begin{minted}[caption={Measured N+1 reproduction and DataLoader collapse (\texttt{nplus1\_demo.py}).},label={lst:ch12-nplus1}]{python}
"""N+1 demonstration: naive per-field store reads vs DataLoader batching.

Runs the same GraphQL operation twice against the same repository and
prints measured store-query counts. Deterministic and offline.

Run:  python nplus1_demo.py
"""

from __future__ import annotations

import asyncio

from catalog_repository import CatalogRepository
from catalog_schema import Context, schema

CATALOG_QUERY = """
query CatalogPage($first: Int!) {
  products(first: $first) {
    edges {
      node {
        name
        category { name }
        reviews { rating }
      }
    }
  }
}
"""


async def run_once(use_loaders: bool) -> tuple[int, int]:
    """Execute CATALOG_QUERY; return (store query count, products returned)."""
    repo = CatalogRepository()
    context = Context(repo, use_loaders=use_loaders)
    result = await schema.execute(
        CATALOG_QUERY, variable_values={"first": 5}, context_value=context
    )
    assert result.errors is None, result.errors
    assert result.data is not None
    return repo.query_count, len(result.data["products"]["edges"])


def main() -> None:
    naive_count, n = run_and_count(use_loaders=False)
    batched_count, _ = run_and_count(use_loaders=True)

    print(f"GraphQL operation asks for {n} products + category + reviews")
    print(f"  naive resolvers   : {naive_count:3d} store queries "
          f"(1 page + {n} category + {n} review reads)")
    print(f"  DataLoader path   : {batched_count:3d} store queries "
          "(1 page + 1 batched category read + 1 batched review read)")
    assert naive_count == 1 + 2 * n
    assert batched_count == 3
    print("N+1 reproduced and collapsed. OK")


def run_and_count(use_loaders: bool) -> tuple[int, int]:
    return asyncio.run(run_once(use_loaders))


if __name__ == "__main__":
    main()
\end{minted}

Eleven reads is $1 + 2N$ at $N=5$; at \texttt{first: 50} it is 101, and
each read is a network round trip against a real database. The DataLoader
path is constant --- 3 reads at any page size --- because batching turns
``one query per product'' into ``one query per \emph{relationship per
tick}.'' The five category reads also demonstrate the cache: products 1
and 2 share category 1, and even the naive count would be worse without
deduplication --- the loader dedupes keys \emph{and} batches them, which
is why the batched count is 3 rather than 5.

\subsection{Listing 4: The Cost-Control Layer}
\label{subsec:ch12-code-cost}

Listing~\ref{lst:ch12-cost} is everything that stands between a curious
client and your database: an AST-walking analyzer that computes depth and
argument-weighted complexity, a budget enforcer, the persisted-query
registry with the server half of the APQ handshake, an AST-based
introspection gate, and \texttt{execute\_graphql\_http} --- the single
entry point that wires them in the order of
Figure~\ref{fig:ch12-pipeline}. Centralizing the pipeline in one function
is the design decision: there is exactly one way for a document to reach
a resolver, so there is exactly one place to audit.

\begin{minted}[caption={Depth limiting, complexity estimation, persisted queries, and the execution gate (\texttt{cost\_control.py}).},label={lst:ch12-cost}]{python}
"""Cost control for GraphQL: depth limiting, complexity estimation, and
persisted queries (allow-list + APQ protocol), plus an HTTP-shaped
execution entry point used by the FastAPI server and the offline tests.

Everything here operates on the *parsed document* before execution, so
expensive operations are rejected without touching a resolver.
"""

from __future__ import annotations

import hashlib
from dataclasses import dataclass
from typing import Any, Awaitable, Callable, Optional

import strawberry
from graphql import (
    DocumentNode,
    FieldNode,
    FragmentDefinitionNode,
    FragmentSpreadNode,
    InlineFragmentNode,
    IntValueNode,
    OperationDefinitionNode,
    parse,
)
from graphql.error import GraphQLSyntaxError

MAX_DEPTH = 7
MAX_COMPLEXITY = 200
DEFAULT_LIST_MULTIPLIER = 10  # assumed page size when `first` is absent


@dataclass(frozen=True)
class OperationMetrics:
    depth: int
    complexity: int


class QueryCostExceeded(Exception):
    """Raised when an operation breaches the depth or complexity budget."""

    def __init__(self, message: str, code: str) -> None:
        super().__init__(message)
        self.code = code


# ---------------------------------------------------------------------------
# Static analysis: depth + complexity
# ---------------------------------------------------------------------------


def _int_argument(field: FieldNode, name: str) -> Optional[int]:
    for arg in field.arguments:
        if arg.name.value == name and isinstance(arg.value, IntValueNode):
            return int(arg.value.value)
    return None


def _list_multiplier(field: FieldNode) -> int:
    """List-returning fields multiply the cost of their sub-selection.

    Only the *top* list field multiplies (the connection); edges/node are
    structural and cost 1, otherwise nesting multiplies twice.
    """
    page = _int_argument(field, "first") or _int_argument(field, "last")
    if page is not None:
        return max(1, min(page, 50))
    if field.name.value in ("products", "reviews"):
        return DEFAULT_LIST_MULTIPLIER
    return 1


def analyze_document(document: DocumentNode) -> OperationMetrics:
    """Compute (max depth, total complexity) over all operations.

    Depth = longest field path, fragment spreads expanded (cycles rejected).
    Complexity = field count, with list fields multiplying their subtree by
    the requested page size (or a conservative default).
    """
    fragments = {
        d.name.value: d
        for d in document.definitions
        if isinstance(d, FragmentDefinitionNode)
    }

    def walk(
        selections: tuple, depth: int, seen: frozenset[str]
    ) -> tuple[int, int]:
        max_depth, cost = depth, 0
        for sel in selections:
            if isinstance(sel, FieldNode):
                child_depth, child_cost = (
                    walk(sel.selection_set.selections, depth + 1, frozenset())
                    if sel.selection_set
                    else (depth + 1, 0)
                )
                cost += 1 + _list_multiplier(sel) * child_cost
                max_depth = max(max_depth, child_depth)
            elif isinstance(sel, InlineFragmentNode):
                d, c = walk(sel.selection_set.selections, depth, seen)
                cost += c
                max_depth = max(max_depth, d)
            elif isinstance(sel, FragmentSpreadNode):
                name = sel.name.value
                if name in seen or name not in fragments:
                    raise QueryCostExceeded(
                        f"invalid fragment spread: {name}", "BAD_FRAGMENT"
                    )
                d, c = walk(
                    fragments[name].selection_set.selections, depth, seen | {name}
                )
                cost += c
                max_depth = max(max_depth, d)
        return max_depth, cost

    worst = OperationMetrics(depth=0, complexity=0)
    for definition in document.definitions:
        if isinstance(definition, OperationDefinitionNode):
            d, c = walk(definition.selection_set.selections, 0, frozenset())
            worst = OperationMetrics(max(worst.depth, d), max(worst.complexity, c))
    return worst


def enforce_cost_budget(
    document: DocumentNode, max_depth: int = MAX_DEPTH, max_cost: int = MAX_COMPLEXITY
) -> OperationMetrics:
    metrics = analyze_document(document)
    if metrics.depth > max_depth:
        raise QueryCostExceeded(
            f"query depth {metrics.depth} exceeds limit {max_depth}",
            "DEPTH_LIMIT_EXCEEDED",
        )
    if metrics.complexity > max_cost:
        raise QueryCostExceeded(
            f"query complexity {metrics.complexity} exceeds limit {max_cost}",
            "COMPLEXITY_LIMIT_EXCEEDED",
        )
    return metrics


# ---------------------------------------------------------------------------
# Persisted queries: allow-list registry + APQ resolution
# ---------------------------------------------------------------------------


class ApqError(Exception):
    def __init__(self, message: str, code: str) -> None:
        super().__init__(message)
        self.code = code


class PersistedQueryRegistry:
    """Allow-list of pre-registered operation documents keyed by SHA-256."""

    def __init__(self) -> None:
        self._by_hash: dict[str, str] = {}

    def register(self, document: str) -> str:
        digest = hashlib.sha256(document.encode("utf-8")).hexdigest()
        self._by_hash[digest] = document
        return digest

    def get(self, digest: str) -> Optional[str]:
        return self._by_hash.get(digest)


def resolve_document(payload: dict[str, Any], registry: PersistedQueryRegistry) -> str:
    """Resolve the operation text from a GraphQL-over-HTTP payload.

    Implements the APQ flow server side: a hash-only payload must already
    be registered; a query+hash payload is verified against its hash and
    then registered.
    """
    query = payload.get("query")
    extensions = payload.get("extensions") or {}
    persisted = extensions.get("persistedQuery") or {}
    digest = persisted.get("sha256Hash")

    if digest is None:
        if not query:
            raise ApqError("no query provided", "BAD_REQUEST")
        return query
    if not isinstance(digest, str) or len(digest) != 64:
        raise ApqError("malformed persistedQuery hash", "BAD_REQUEST")
    if query is None:
        stored = registry.get(digest)
        if stored is None:
            raise ApqError("persisted query not found", "PersistedQueryNotFound")
        return stored
    actual = hashlib.sha256(query.encode("utf-8")).hexdigest()
    if actual != digest:
        raise ApqError("provided hash does not match query", "PersistedQueryMismatch")
    registry.register(query)
    return query


def uses_introspection(document: DocumentNode) -> bool:
    """True if any selected field is a meta field (__schema / __type).

    AST-based, not substring-based: aliases and fragments cannot hide it.
    (``__typename`` is excluded: it is cheap and used by client caches.)
    """
    def walk(selections: tuple) -> bool:
        for sel in selections:
            if isinstance(sel, FieldNode):
                if sel.name.value in ("__schema", "__type"):
                    return True
                if sel.selection_set and walk(sel.selection_set.selections):
                    return True
            elif isinstance(sel, (InlineFragmentNode, FragmentDefinitionNode)):
                if walk(sel.selection_set.selections):
                    return True
        return False

    for definition in document.definitions:
        if hasattr(definition, "selection_set"):
            if walk(tuple(definition.selection_set.selections)):
                return True
    return False


# ---------------------------------------------------------------------------
# HTTP-shaped execution entry point (used by server.py and tests)
# ---------------------------------------------------------------------------


def _error_body(message: str, code: str) -> dict[str, Any]:
    return {"errors": [{"message": message, "extensions": {"code": code}}]}


async def execute_graphql_http(
    payload: dict[str, Any],
    *,
    schema: strawberry.Schema,
    registry: PersistedQueryRegistry,
    context_factory: Callable[[], Any],
    max_depth: int = MAX_DEPTH,
    max_cost: int = MAX_COMPLEXITY,
    allow_introspection: bool = False,
) -> tuple[int, dict[str, Any]]:
    """Parse -> resolve persisted -> budget-check -> execute.

    Returns (HTTP status, response body) with GraphQL semantics: execution
    errors are HTTP 200 with a populated ``errors`` array; malformed or
    over-budget requests are HTTP 400.
    """
    try:
        document_text = resolve_document(payload, registry)
    except ApqError as exc:
        # APQ misses are HTTP 200 per the protocol: the client retries.
        status = 200 if exc.code.startswith("PersistedQuery") else 400
        return status, _error_body(str(exc), exc.code)

    try:
        document = parse(document_text)
    except GraphQLSyntaxError as exc:
        return 400, _error_body(f"syntax error: {exc.message}", "GRAPHQL_PARSE_FAILED")

    if not allow_introspection and uses_introspection(document):
        return 400, _error_body(
            "introspection is disabled", "INTROSPECTION_DISABLED"
        )

    try:
        enforce_cost_budget(document, max_depth, max_cost)
    except QueryCostExceeded as exc:
        return 400, _error_body(str(exc), exc.code)

    result = await schema.execute(
        document_text,
        variable_values=payload.get("variables"),
        operation_name=payload.get("operationName"),
        context_value=context_factory(),
    )
    body: dict[str, Any] = {"data": result.data}
    if result.errors:
        body["errors"] = [err.formatted for err in result.errors]
    return 200, body
\end{minted}

The complexity estimator encodes one deliberate modeling choice:
\texttt{\_list\_multiplier} charges the \emph{connection} field
(\texttt{products}) for its requested page size but treats the structural
\texttt{edges}/\texttt{node} wrappers as cost 1 --- otherwise nesting
multiplies twice and a modest \texttt{first: 2} query reads as cost 441.
Cost models are heuristics with adversaries; calibrate them against real
traffic and keep the timeout backstop armed.

\subsection{Listing 5: The FastAPI Host}
\label{subsec:ch12-code-server}

The server (Listing~\ref{lst:ch12-server}) is thin by design: parse JSON,
call \texttt{execute\_graphql\_http}, return. Introspection policy is an
environment flag defaulting to \emph{off}; the registry is seeded with the
mobile app's catalog operation so allow-listed traffic never pays the APQ
round trip. There is no GraphQL router magic here on purpose --- when the
endpoint is thirty lines you own, the security review is a short meeting.

\begin{minted}[caption={FastAPI host with introspection gating and a seeded allow-list (\texttt{server.py}).},label={lst:ch12-server}]{python}
"""FastAPI host for the catalog schema.

The HTTP layer is deliberately thin: all GraphQL policy (persisted-query
resolution, depth/complexity budgets, execution) lives in cost_control.py
so it is testable without a web framework. Introspection is gated by an
explicit policy flag, never by "we forgot".

Run (dev only):  uvicorn server:app --port 8000
"""

from __future__ import annotations

import os
from typing import Any

from fastapi import FastAPI, Request
from fastapi.responses import JSONResponse

from catalog_repository import CatalogRepository
from catalog_schema import Context, schema
from cost_control import PersistedQueryRegistry, execute_graphql_http

INTROSPECTION_ENABLED = os.environ.get("NWR_INTROSPECTION", "off") == "on"

_registry = PersistedQueryRegistry()
# Pre-register the operations the mobile app ships with (allow-list seed).
_CATALOG_OP = """
query CatalogPage($first: Int!) {
  products(first: $first) {
    edges { node { name category { name } } }
    pageInfo { hasNextPage endCursor }
  }
}
"""
CATALOG_OP_HASH = _registry.register(_CATALOG_OP)


def _context_factory() -> Context:
    # One repository + one set of DataLoaders per request.
    return Context(CatalogRepository())


def create_app() -> FastAPI:
    app = FastAPI(title="Northwind Robotics Catalog", version="1.0.0")

    @app.post("/graphql")
    async def graphql_endpoint(request: Request) -> JSONResponse:
        payload: dict[str, Any] = await request.json()
        status, body = await execute_graphql_http(
            payload,
            schema=schema,
            registry=_registry,
            context_factory=_context_factory,
            allow_introspection=INTROSPECTION_ENABLED,
        )
        return JSONResponse(body, status_code=status)

    @app.get("/health")
    async def health() -> dict[str, str]:
        return {"status": "ok"}

    return app


app = create_app()
\end{minted}

\subsection{Listing 6: The httpx Client}
\label{subsec:ch12-code-client}

Listing~\ref{lst:ch12-client} is the client half of the contract. It
separates the three failure channels the wire shape creates:
\texttt{GraphQLTransportError} for transport failure and 5xx,
\texttt{GraphQLFieldErrors} for requests that never executed (4xx) and
for callers who prefer exceptions, and a returned
\texttt{GraphQLResponse} carrying \texttt{data} \emph{and}
\texttt{errors} together for the partial-success case --- the client
library does not get to decide whether partial data is usable; the
application does. \texttt{execute\_persisted} implements the APQ
handshake: hash first, retry with full text only on
\texttt{PersistedQueryNotFound}. The demo runs entirely in-process through
\texttt{httpx.ASGITransport} \cite{httpxdocs}: real HTTP semantics, zero
sockets.

\begin{minted}[caption={GraphQL-over-HTTP client with variables, error channels, and APQ (\texttt{graphql\_client.py}).},label={lst:ch12-client}]{python}
"""GraphQL-over-HTTP client with variables, error handling, and APQ.

Transport is injected, so the demo below runs fully offline against the
in-process FastAPI app (httpx.ASGITransport) -- no socket is opened.

Run:  python graphql_client.py
"""

from __future__ import annotations

import asyncio
import hashlib
import json
from dataclasses import dataclass
from typing import Any, Optional

import httpx

GRAPHQL_PATH = "/graphql"


class GraphQLTransportError(Exception):
    """Network failure or non-GraphQL HTTP error (5xx, timeouts)."""


class GraphQLFieldErrors(Exception):
    """The response carried a populated ``errors`` array (HTTP 200)."""

    def __init__(self, errors: list[dict[str, Any]], data: Any) -> None:
        messages = [e.get("message", "?") for e in errors]
        super().__init__("; ".join(messages))
        self.errors = errors
        self.data = data  # partial data may still be present


@dataclass(frozen=True)
class GraphQLResponse:
    data: Any
    errors: list[dict[str, Any]]


class GraphQLClient:
    def __init__(self, base_url: str, transport: Optional[httpx.AsyncBaseTransport] = None,
                 timeout: float = 10.0) -> None:
        self._client = httpx.AsyncClient(
            base_url=base_url, transport=transport, timeout=timeout
        )

    async def close(self) -> None:
        await self._client.aclose()

    async def __aenter__(self) -> "GraphQLClient":
        return self

    async def __aexit__(self, *exc: object) -> None:
        await self.close()

    async def execute(
        self,
        query: str,
        variables: Optional[dict[str, Any]] = None,
        operation_name: Optional[str] = None,
    ) -> GraphQLResponse:
        payload: dict[str, Any] = {"query": query}
        if variables is not None:
            payload["variables"] = variables
        if operation_name is not None:
            payload["operationName"] = operation_name
        return await self._post(payload)

    async def execute_persisted(
        self, query: str, variables: Optional[dict[str, Any]] = None
    ) -> GraphQLResponse:
        """Automatic Persisted Queries: hash first, retry with full text
        only when the server answers PersistedQueryNotFound."""
        digest = hashlib.sha256(query.encode("utf-8")).hexdigest()
        ext = {"persistedQuery": {"version": 1, "sha256Hash": digest}}
        first = await self._post({"extensions": ext, "variables": variables or {}})
        if not self._is_apq_miss(first):
            return first
        # Retry once, sending query text + hash so the server can register it.
        return await self._post(
            {"query": query, "extensions": ext, "variables": variables or {}}
        )

    @staticmethod
    def _is_apq_miss(response: GraphQLResponse) -> bool:
        return any(
            (e.get("extensions") or {}).get("code") == "PersistedQueryNotFound"
            for e in response.errors
        )

    async def _post(self, payload: dict[str, Any]) -> GraphQLResponse:
        try:
            response = await self._client.post(
                GRAPHQL_PATH,
                content=json.dumps(payload),
                headers={"Content-Type": "application/json"},
            )
        except httpx.HTTPError as exc:
            raise GraphQLTransportError(f"transport failure: {exc}") from exc
        if response.status_code >= 500:
            raise GraphQLTransportError(f"server error: HTTP {response.status_code}")
        body = response.json()
        errors = body.get("errors") or []
        # 4xx here means the request never executed (parse/cost rejection).
        if response.status_code >= 400:
            raise GraphQLFieldErrors(errors, data=None)
        return GraphQLResponse(data=body.get("data"), errors=errors)


CATALOG_QUERY = """
query CatalogPage($first: Int!) {
  products(first: $first) {
    edges { node { name category { name } } }
    pageInfo { hasNextPage endCursor }
  }
}
"""


async def demo() -> None:
    import logging

    from server import create_app  # in-process app: zero network

    # Strawberry logs every field error; keep the demo output readable.
    logging.getLogger("strawberry.execution").setLevel(logging.CRITICAL)

    transport = httpx.ASGITransport(app=create_app())
    async with GraphQLClient("http://testserver", transport=transport) as client:
        # 1. Plain query with variables.
        page = await client.execute(CATALOG_QUERY, variables={"first": 2})
        edges = page.data["products"]["edges"]
        print(f"plain query   : {len(edges)} products, "
              f"hasNextPage={page.data['products']['pageInfo']['hasNextPage']}")

        # 2. APQ: first call misses, second is served hash-only.
        await client.execute_persisted(CATALOG_QUERY, variables={"first": 2})
        cached = await client.execute_persisted(CATALOG_QUERY, variables={"first": 3})
        print(f"APQ round trip: {len(cached.data['products']['edges'])} products")

        # 3. Field-level error: partial data + errors, HTTP 200.
        bad = await client.execute(
            "query { product(productId: 4) { name marginPercent } }"
        )
        print(f"field error   : data={bad.data} "
              f"code={bad.errors[0]['extensions']['code']}")

        # 4. Rejected request: over-depth query is refused pre-execution.
        try:
            deep = (
                "query Deep { products { edges { node { "
                + "reviews { " * 12 + "rating" + " }" * 12
                + " } } } }"
            )
            await client.execute(deep)
        except GraphQLFieldErrors as exc:
            print(f"depth rejected: {exc.errors[0]['extensions']['code']}")


if __name__ == "__main__":
    asyncio.run(demo())
\end{minted}

Run it and you get the whole pipeline in four lines of output --- a
paginated query, an APQ round trip (miss, register, hash-only hit), a
field error arriving as partial data, and a depth rejection:

\begin{minted}{text}
plain query   : 2 products, hasNextPage=True
APQ round trip: 3 products
field error   : data={'product': None} code=COST_DATA_UNAVAILABLE
depth rejected: DEPTH_LIMIT_EXCEEDED
\end{minted}

\subsection{Listing 7: The Test Suite}
\label{subsec:ch12-code-tests}

Listing~\ref{lst:ch12-tests} is 21 tests, all in-process via
\texttt{schema.execute} or \texttt{execute\_graphql\_http}, driven by
\texttt{asyncio.run} so no pytest plugins are required. It covers the
contract surface a reviewer should demand: connection walking without
gaps or duplicates, cursor forgery, aliases and \texttt{@skip}, both
halves of null propagation (non-null error voids the product; the
nullable list element voids only itself), all four typed mutation errors,
the $N{+}1$ counts as \emph{assertions} (11 naive, 3 batched), depth and
complexity rejection, the full APQ handshake including hash mismatch, and
introspection denied by default and served when explicitly gated on.

\begin{minted}[caption={In-process test suite: 21 tests, zero network (\texttt{test\_catalog.py}).},label={lst:ch12-tests}]{python}
"""In-process test suite for the Northwind Robotics GraphQL catalog.

No sockets, no network, deterministic seed data. Async execution is driven
with asyncio.run so no pytest plugins are required.

Run:  pytest test_catalog.py -q
"""

from __future__ import annotations

import asyncio
import logging
from typing import Any, Optional

import pytest

from catalog_repository import CatalogRepository
from catalog_schema import Context, schema
from cost_control import (
    PersistedQueryRegistry,
    analyze_document,
    enforce_cost_budget,
    QueryCostExceeded,
    execute_graphql_http,
)
from graphql import parse

logging.getLogger("strawberry.execution").setLevel(logging.CRITICAL)


def run(
    query: str,
    variables: Optional[dict[str, Any]] = None,
    repo: Optional[CatalogRepository] = None,
):
    """Execute against a fresh repository unless one is supplied."""
    repo = repo or CatalogRepository()
    context = Context(repo)
    result = asyncio.run(
        schema.execute(query, variable_values=variables, context_value=context)
    )
    return result, repo


# ---------------------------------------------------------------------------
# Queries, variables, aliases, fragments
# ---------------------------------------------------------------------------


def test_products_connection_first_page() -> None:
    result, _ = run(
        """
        query Page($first: Int!) {
          products(first: $first) {
            edges { cursor node { sku name } }
            pageInfo { hasNextPage hasPreviousPage startCursor endCursor }
            totalCount
          }
        }
        """,
        variables={"first": 2},
    )
    assert result.errors is None
    data = result.data["products"]
    assert [e["node"]["sku"] for e in data["edges"]] == ["NWR-ACT-100", "NWR-ACT-200"]
    assert data["pageInfo"]["hasNextPage"] is True
    assert data["pageInfo"]["hasPreviousPage"] is False
    assert data["totalCount"] == 5


def test_connection_walk_is_gapless() -> None:
    """Walking by cursor visits every product exactly once."""
    seen: list[str] = []
    after: Optional[str] = None
    for _ in range(10):
        result, _ = run(
            """
            query Walk($first: Int!, $after: String) {
              products(first: $first, after: $after) {
                edges { cursor node { sku } }
                pageInfo { hasNextPage endCursor }
              }
            }
            """,
            variables={"first": 2, "after": after},
        )
        assert result.errors is None
        page = result.data["products"]
        seen.extend(e["node"]["sku"] for e in page["edges"])
        if not page["pageInfo"]["hasNextPage"]:
            break
        after = page["pageInfo"]["endCursor"]
    assert seen == [
        "NWR-ACT-100", "NWR-ACT-200", "NWR-SEN-310", "NWR-SEN-410", "NWR-CTL-900",
    ]


def test_invalid_cursor_returns_field_error() -> None:
    result, _ = run(
        'query { products(first: 2, after: "%%%") { edges { cursor } } }'
    )
    assert result.data == {"products": None}
    assert result.errors is not None
    assert result.errors[0].extensions["code"] == "BAD_CURSOR"


def test_status_filter_and_alias() -> None:
    result, _ = run(
        """
        query {
          active: products(first: 10, status: ACTIVE) { edges { node { sku } } }
          retired: products(first: 10, status: DISCONTINUED) {
            edges { node { sku } }
          }
        }
        """
    )
    assert result.errors is None
    assert [e["node"]["sku"] for e in result.data["active"]["edges"]] == [
        "NWR-ACT-100", "NWR-ACT-200", "NWR-SEN-310", "NWR-CTL-900",
    ]
    assert [e["node"]["sku"] for e in result.data["retired"]["edges"]] == [
        "NWR-SEN-410"
    ]


def test_fragments_and_skip_directive() -> None:
    result, _ = run(
        """
        query F($withReviews: Boolean!) {
          product(productId: 1) {
            ...ProductCard
          }
        }
        fragment ProductCard on Product {
          name
          reviews @skip(if: $withReviews) { rating }
        }
        """,
        variables={"withReviews": True},
    )
    assert result.errors is None
    assert result.data == {"product": {"name": "Servo Linear Actuator 100mm"}}


# ---------------------------------------------------------------------------
# Nullability propagation on errors
# ---------------------------------------------------------------------------


def test_non_null_field_error_nulls_the_product() -> None:
    """marginPercent is Int! and raises for product 4 (no cost data).

    Per the spec's null-propagation rule the error bubbles to the nearest
    nullable position: the whole `product` field becomes null, and the
    response is partial: HTTP 200 with data + errors.
    """
    result, _ = run("query { product(productId: 4) { name marginPercent } }")
    assert result.data == {"product": None}
    assert result.errors is not None
    error = result.errors[0]
    assert error.path == ["product", "marginPercent"]
    assert error.extensions["code"] == "COST_DATA_UNAVAILABLE"


def test_nullable_sibling_survives_in_list() -> None:
    """Inside a nullable list element, only that element nulls out."""
    result, _ = run(
        """
        query {
          products(first: 5) {
            edges { node { sku marginPercent } }
          }
        }
        """
    )
    assert result.errors is not None
    nodes = [e["node"] for e in result.data["products"]["edges"]]
    assert nodes[3] is None  # product 4 nulled by propagation
    assert nodes[0]["marginPercent"] == 54  # siblings intact


# ---------------------------------------------------------------------------
# Mutations: typed errors via union payloads
# ---------------------------------------------------------------------------


ADD_REVIEW = """
mutation Add($input: AddReviewInput!) {
  addReview(input: $input) {
    __typename
    ... on AddReviewSuccess { review { id rating headline } }
    ... on AddReviewError { code message }
  }
}
"""


def test_add_review_success() -> None:
    result, _ = run(
        ADD_REVIEW,
        variables={"input": {"productId": 2, "rating": 4, "headline": "Solid"}},
    )
    assert result.errors is None
    payload = result.data["addReview"]
    assert payload["__typename"] == "AddReviewSuccess"
    assert payload["review"]["rating"] == 4


@pytest.mark.parametrize(
    "input_vars, code",
    [
        ({"productId": 2, "rating": 0, "headline": "Bad"}, "RATING_OUT_OF_RANGE"),
        ({"productId": 2, "rating": 6, "headline": "Bad"}, "RATING_OUT_OF_RANGE"),
        ({"productId": 2, "rating": 5, "headline": "   "}, "EMPTY_HEADLINE"),
        ({"productId": 999, "rating": 5, "headline": "Ghost"}, "PRODUCT_NOT_FOUND"),
    ],
)
def test_add_review_typed_errors(input_vars: dict[str, Any], code: str) -> None:
    result, _ = run(ADD_REVIEW, variables={"input": input_vars})
    assert result.errors is None  # expected failure is data, not an exception
    payload = result.data["addReview"]
    assert payload["__typename"] == "AddReviewError"
    assert payload["code"] == code


# ---------------------------------------------------------------------------
# N+1: measured query counts
# ---------------------------------------------------------------------------


CATALOG_QUERY = """
query CatalogPage($first: Int!) {
  products(first: $first) {
    edges { node { name category { name } reviews { rating } } }
  }
}
"""


def _count_queries(use_loaders: bool) -> int:
    repo = CatalogRepository()
    context = Context(repo, use_loaders=use_loaders)
    result = asyncio.run(
        schema.execute(
            CATALOG_QUERY, variable_values={"first": 5}, context_value=context
        )
    )
    assert result.errors is None
    return repo.query_count


def test_nplus1_naive_explosion() -> None:
    assert _count_queries(use_loaders=False) == 11  # 1 + 5 + 5


def test_dataloader_collapses_nplus1() -> None:
    assert _count_queries(use_loaders=True) == 3  # 1 + 1 + 1


# ---------------------------------------------------------------------------
# Cost control: depth, complexity, persisted queries
# ---------------------------------------------------------------------------


def test_depth_limiter_rejects_deep_query() -> None:
    deep = (
        "query Deep { products { edges { node { "
        + "reviews { " * 12 + "rating" + " }" * 12 + " } } } }"
    )
    with pytest.raises(QueryCostExceeded) as excinfo:
        enforce_cost_budget(parse(deep))
    assert excinfo.value.code == "DEPTH_LIMIT_EXCEEDED"


def test_complexity_scales_with_page_size() -> None:
    small = analyze_document(
        parse("query { products(first: 2) { edges { node { name } } } }")
    )
    large = analyze_document(
        parse("query { products(first: 50) { edges { node { name } } } }")
    )
    assert large.complexity > small.complexity


def test_complexity_limit_rejects_god_query() -> None:
    god = "query { " + " ".join(
        f"p{i}: products(first: 50) {{ edges {{ node {{ name reviews {{ rating }} }} }} }}"
        for i in range(8)
    ) + " }"
    with pytest.raises(QueryCostExceeded) as excinfo:
        enforce_cost_budget(parse(god))
    assert excinfo.value.code == "COMPLEXITY_LIMIT_EXCEEDED"


def test_persisted_query_apq_flow() -> None:
    registry = PersistedQueryRegistry()
    factory = lambda: Context(CatalogRepository())  # noqa: E731

    async def apq_round_trip() -> tuple[int, dict[str, Any]]:
        import hashlib

        digest = hashlib.sha256(CATALOG_QUERY.encode("utf-8")).hexdigest()
        ext = {"persistedQuery": {"version": 1, "sha256Hash": digest}}
        # Step 1: hash only -> PersistedQueryNotFound (HTTP 200)
        status, body = await execute_graphql_http(
            {"extensions": ext, "variables": {"first": 2}},
            schema=schema, registry=registry, context_factory=factory,
        )
        assert status == 200
        assert body["errors"][0]["extensions"]["code"] == "PersistedQueryNotFound"
        # Step 2: query + hash -> executes and registers
        return await execute_graphql_http(
            {"query": CATALOG_QUERY, "extensions": ext, "variables": {"first": 2}},
            schema=schema, registry=registry, context_factory=factory,
        )

    status, body = asyncio.run(apq_round_trip())
    assert status == 200
    assert "errors" not in body
    assert len(body["data"]["products"]["edges"]) == 2


def test_persisted_query_hash_mismatch_rejected() -> None:
    registry = PersistedQueryRegistry()

    async def mismatch() -> tuple[int, dict[str, Any]]:
        return await execute_graphql_http(
            {
                "query": "query { product(productId: 1) { name } }",
                "extensions": {"persistedQuery": {"sha256Hash": "0" * 64}},
            },
            schema=schema,
            registry=registry,
            context_factory=lambda: Context(CatalogRepository()),
        )

    status, body = asyncio.run(mismatch())
    assert status == 200
    assert body["errors"][0]["extensions"]["code"] == "PersistedQueryMismatch"


def test_introspection_blocked_by_default() -> None:
    async def probe() -> tuple[int, dict[str, Any]]:
        return await execute_graphql_http(
            {"query": "{ __schema { types { name } } }"},
            schema=schema,
            registry=PersistedQueryRegistry(),
            context_factory=lambda: Context(CatalogRepository()),
        )

    status, body = asyncio.run(probe())
    assert status == 400
    assert body["errors"][0]["extensions"]["code"] == "INTROSPECTION_DISABLED"


def test_introspection_allowed_when_gated_on() -> None:
    async def probe() -> tuple[int, dict[str, Any]]:
        return await execute_graphql_http(
            {"query": "{ __schema { queryType { name } } }"},
            schema=schema,
            registry=PersistedQueryRegistry(),
            context_factory=lambda: Context(CatalogRepository()),
            allow_introspection=True,
        )

    status, body = asyncio.run(probe())
    assert status == 200
    assert body["data"]["__schema"]["queryType"]["name"] == "Query"
\end{minted}

\begin{minted}{text}
.....................                                                    [100%]
21 passed in 2.78s
\end{minted}

%% ============================================================================
\section{Common Pitfalls \& Anti-Patterns}
\label{sec:ch12-pitfalls}

Each entry names the failure mode, why it happens, and the smallest
durable fix. Most are defaults --- GraphQL ships insecure-by-default in
the specific sense that its ergonomics and its guardrails are both
opt-in.

\subsection{N+1 by Construction}
\textbf{Failure.} Every relationship field performs its own store read;
latency and database load scale with page size, not with request count.
\textbf{Why.} Field-level resolution is the default physics; nothing in a
naive resolver hints it is one of fifty identical calls.
\textbf{Fix.} DataLoader per request, batch reads in the repository,
\emph{count queries in tests} --- our suite asserts 3, not ``fast.''

\subsection{Open Introspection in Production}
\textbf{Failure.} \texttt{\_\_schema} answers anyone who asks with your
complete type inventory. \textbf{Why.} Frameworks enable it by default
because GraphiQL needs it. \textbf{Fix.} Gate it at the validation layer
(AST check, not substring), default off in production, on in development;
treat the schema as an asset with an audience.

\subsection{Unlimited Depth and Complexity}
\textbf{Failure.} A cyclic traversal or a \texttt{first: 10000} god-query
pins the executor and the database. \textbf{Why.} The client chooses the
work; unguarded servers accept all of it. \textbf{Fix.} Static depth and
argument-weighted complexity budgets between parse and execution, plus
resolver timeouts for estimation errors.

\subsection{Treating GraphQL Errors as HTTP 500 Only}
\textbf{Failure.} The client checks \texttt{status\_code}, sees 200,
renders \texttt{data.product.name} --- which is \texttt{null} --- and the
stack trace ships to a user. Or the inverse: any \texttt{errors} entry
discards a usable partial page. \textbf{Why.} Two error channels with
different semantics, and REST instincts cover only one. \textbf{Fix.}
Client libraries expose transport, field-error, and partial-data as
distinct outcomes (Listing~\ref{lst:ch12-client}); clients render partial
data where the schema's nullability says it is safe.

\subsection{Schema Equals Database Schema}
\textbf{Failure.} Tables are exported verbatim as types; every
re-normalization is a breaking API change; internal columns become public
surface. \textbf{Why.} Code generation from the ORM is ten minutes of
work. \textbf{Fix.} Design the schema as the \emph{product contract}
--- client-shaped, not table-shaped --- and let resolvers be the
anti-corruption layer. Our \texttt{Product} hides
\texttt{category\_id} (a storage detail) behind \texttt{category} (a
concept).

\subsection{God-Queries from Your Own Clients}
\textbf{Failure.} One team's dashboard composes a 40-field, depth-9
operation that runs every 30 seconds from every open browser; the
incident review blames ``traffic.'' \textbf{Why.} Composition is the
feature; nobody budgets it. \textbf{Fix.} Persisted queries for first-party
clients (the set of operations is known at build time), per-client cost
budgets, and complexity telemetry per operation name so the expensive
ones have owners.

\subsection{DataLoader Shared Across Requests}
\textbf{Failure.} A module-level loader caches user A's record and serves
it to user B. \textbf{Why.} The cache is the feature; its scope is the
footgun --- it must die with the request. \textbf{Fix.} Construct loaders
in the context factory, one set per request, as
Listing~\ref{lst:ch12-schema} does; never in module scope.

\subsection{Bypassing the HTTP Cache Without a Plan}
\textbf{Failure.} Everything is POST \texttt{/graphql}; the CDN is a
pass-through; a product launch melts the origin that REST would have
served from cache. \textbf{Why.} The default GraphQL-over-HTTP mapping
discards the URL-as-cache-key. \textbf{Fix.} Decide deliberately: APQ
over GET restores cache keys for read-heavy public traffic; private,
per-user data was never CDN-cacheable anyway --- but know which you are
serving before you choose the transport \cite{rfc9110}.

%% ============================================================================
\section{Hands-On Production Lab \& Real-World Case Study}
\label{sec:ch12-lab}

\subsection{Lab: Build, Break, and Harden the Catalog API}

Four hours, one Windows workstation, zero network. Work in
\texttt{C:\textbackslash{}training\textbackslash{}graphql\_lab}.

\begin{enumerate}
  \item \textbf{Stand up the lab.} Create the six files of
        Section~\ref{sec:ch12-code}. Then:
\begin{minted}{text}
python -m venv .venv
.\.venv\Scripts\Activate.ps1
pip install "strawberry-graphql[fastapi]>=0.230" "fastapi>=0.110" "httpx>=0.27" "pytest>=8"
pytest -q
\end{minted}
        All 21 tests pass. You now hold a working GraphQL service with its
        guardrails already attached --- read them before you remove them.
  \item \textbf{Reproduce the N+1.} Run \texttt{python nplus1\_demo.py}
        and record the counts (11 naive, 3 batched). Then edit the schema
        so \texttt{use\_loaders} defaults to \texttt{False} and re-run
        \texttt{pytest}: \texttt{test\_dataloader\_collapses\_nplus1}
        fails. Revert. You have now seen the failure as a test assertion,
        which is the only way it will ever stay fixed.
  \item \textbf{Find the threshold.} In a Python REPL, call
        \texttt{analyze\_document(parse(...))} on
        \texttt{CATALOG\_QUERY} with \texttt{first} $= 2, 10, 50$ and
        record complexity. Then raise it until
        \texttt{COMPLEXITY\_LIMIT\_EXCEEDED} fires. Note the page size at
        which a \emph{legitimate} query dies --- that number is a product
        decision, not an implementation detail.
  \item \textbf{Break the allow-list.} Using the client, POST a hash-only
        request for an unregistered operation and observe
        \texttt{PersistedQueryNotFound}; then send a mismatched hash/text
        pair and observe \texttt{PersistedQueryMismatch}. Explain in one
        sentence why the mismatch check exists (hint: the hash is the
        allow-list key).
  \item \textbf{Harden the endpoint.} With introspection off, attempt
        \texttt{\{ \_\_schema \{ types \{ name \} \} \} } and an
        aliased variant; both must fail with
        \texttt{INTROSPECTION\_DISABLED}. Then try the array-of-queries
        batch attack (\texttt{[\{"query": ...\}, \{"query": ...\}, ...]})
        against \texttt{server.py}: it is rejected (the endpoint accepts
        only objects) --- state which control catches it and what you
        would add if batching were a product requirement.
  \item \textbf{Write the regression.} Add a test that walks the entire
        connection with \texttt{first: 2} and asserts the five SKUs in
        order with no duplicates --- compare with
        \texttt{test\_connection\_walk\_is\_gapless}, which already does
        this. If yours is identical, write a different one: cursor reuse
        after a mutation that adds a review must not shift the product
        ordering.
\end{enumerate}

\subsection{Case Study: The Introspection-Enabled Enumeration Incident}
\label{sec:ch12-casestudy}

\emph{Composite reconstruction from public incident patterns; all names,
figures, and timelines fictionalized.}

A mid-size robotics-parts marketplace --- call it Northwind Robotics'
commerce arm --- shipped a GraphQL gateway in front of its catalog,
pricing, and dealer-account services to feed a new mobile app. Launch
went well. Nine weeks later, a security researcher's responsible
disclosure landed: a complete map of every type and field in the graph,
\emph{including two that never appeared in any shipped client} ---
\texttt{dealerAccount \{ terms \{ creditLimitCents \} \}} and an
undocumented \texttt{internalNotes} field on \texttt{Product}.

The disclosure's reproduction steps were four lines long, and that is the
lesson. Step one: introspection was enabled in production because nobody
had turned it off --- GraphiQL had been ``temporarily'' exposed for the
mobile team and the temporary flag ossified. Step two: with the field
inventory in hand, the researcher composed valid queries against
\texttt{dealerAccount} in minutes; field-level authorization existed on
\emph{most} account fields, but \texttt{creditLimitCents} had been added
in a hurry and its resolver checked only authentication, not
dealer-scoping --- any authenticated dealer could page through every
dealer's credit terms, with cursor pagination doing what it is designed
to do: make enumeration complete and cheap. Step three: the gateway's
only rate limit counted \emph{requests}; the enumeration ran inside
array-of-queries batches of fifty, so a full exfiltration of 38{,}000
dealer records presented to the WAF as 760 ordinary-looking POSTs over a
weekend. Step four: the cost gate did not exist. Depth and complexity
limits had been ticketed as ``post-launch hardening.''

The remediation reads like this chapter's checklist, because this chapter
is written from incidents like this one. Introspection and GraphiQL gated
behind the same explicit production flag, default off. Authorization
moved into a resolver decorator with a deny-by-default audit --- every
field returning account data declares its required scope, and CI fails on
undeclared fields. Array batching capped at five operations with cost
summed across the batch. Rate limiting re-keyed from requests to
estimated complexity units per token per minute. And the two shadow
fields were removed from the schema, because the deepest finding was not
a bug but a governance failure: the schema had drifted from the product
contract, and introspection made the drift public. The total engineering
cost of the remediation was about three engineer-weeks; the incident
review estimated the enumeration itself at under \$50 of compute ---
attackers get the bulk discount that your missing cost controls hand
them.

\begin{warningbox}
Every control in the remediation was a \emph{default} somewhere waiting to
be flipped: introspection on, batching unbounded, request-count rate
limits, authorization-by-convention. A production GraphQL review is
substantially a review of defaults. Write them down as explicit policy ---
the lab's step 5 --- or the incident review will write them down for you.
\end{warningbox}

%% ============================================================================
\section{Self-Assessment Exercises \& Deep-Dive Questions}
\label{sec:ch12-exercises}

\subsection{Recall and Comprehension}

\begin{enumerate}
  \item List the eight kinds in the GraphQL type system and give a
        catalog-domain example of each.
  \item Write the four list/nullability combinations for
        \texttt{[Review]} and state, for each, what an error in one
        element does to the response.
  \item What are the three phases of request processing, and which
        failure channel does each use?
  \item Why can variables not be used to inject into the query document?
        What injection remains possible?
  \item In the APQ handshake, what does the client send first, what does
        the server answer on a miss, and why is the miss an HTTP 200?
\end{enumerate}

\subsection{Application}

\begin{enumerate}[resume]
  \item Extend \texttt{catalog\_schema.py} with a
        \texttt{Category.products} field. Predict the $N{+}1$ count for a
        query selecting all categories and their products' names before
        you run it; then instrument and verify.
  \item Our estimator charges connection fields by their \texttt{first}
        argument. Construct a query that is cheap by our estimator but
        expensive to execute. What estimator change closes the hole, and
        what legitimate traffic does it harm?
  \item Implement \texttt{last}/\texttt{before} backward pagination in
        the repository and schema. What does \texttt{hasPreviousPage}
        require that forward pagination did not?
  \item Add a \texttt{@deprecated} field to the schema (e.g.\ rename
        \texttt{unitPriceCents} to \texttt{priceCents} and deprecate the
        old one). Which client-facing guarantee of continuous evolution
        does this exercise?
  \item Using only Listing~\ref{lst:ch12-client}'s error channels, write
        client logic that renders a product card from partial data when
        \texttt{COST\_DATA\_UNAVAILABLE} is the only error, but retries
        the operation on \texttt{GraphQLTransportError}.
\end{enumerate}

\subsection{Deep-Dive Questions}

\begin{enumerate}[resume]
  \item The specification executes mutation root fields serially but query
        fields without ordering. Why is serial mutation execution the only
        ordering guarantee the spec can make, and what does it imply for
        client design?
  \item Null propagation couples error containment to schema design.
        Argue for and against declaring \texttt{edges: [ProductEdge!]!} in
        a public schema. What telemetry would tell you whether you chose
        correctly?
  \item Persisted-query allow-lists turn an open query language into a
        closed RPC surface. What does GraphQL still buy you under a full
        allow-list, and what does it cost?
  \item Our complexity model multiplies by page size at the connection
        only. Sketch a model that also prices \emph{resolver heterogeneity}
        (a field hitting a search index costs more than a field hitting
        memory). Where does per-field cost metadata live in a code-first
        schema?
  \item Federation moves the $N{+}1$ problem from resolver-to-database to
        gateway-to-subgraph. Does DataLoader-style batching transfer to
        entity resolution across services? What breaks?
\end{enumerate}

%% ============================================================================
\section{Interview Q\&A Vault}
\label{sec:ch12-interview-vault}

Ordered junior $\to$ staff. Answers are model responses at the depth a
strong candidate gives verbally.

\begin{enumerate}[label=\textbf{Q\arabic*.}, leftmargin=2.6em]
  % ---- junior ----
  \item \textbf{What is GraphQL, in one paragraph?}
        A query language and runtime for APIs: the server publishes a
        typed schema, the client sends an operation selecting exactly the
        fields it needs, and the response mirrors that selection. One
        endpoint, client-composed shapes, strong typing, introspection.
  \item \textbf{How is a GraphQL request different from a REST request?}
        REST addresses a resource whose shape the server fixed; GraphQL
        POSTs a document describing the desired shape to one endpoint.
        REST spreads a screen's data across $k$ round trips; GraphQL
        composes it into one --- at the price of server-side cost
        governance.
  \item \textbf{Name the built-in scalar types.}
        \texttt{Int}, \texttt{Float}, \texttt{String}, \texttt{Boolean},
        \texttt{ID}. \texttt{ID} serializes as a string and means
        ``opaque; do not parse.'' Custom scalars (e.g.\
        \texttt{DateTime}) define their own serialization.
  \item \textbf{What is the difference between an interface and a union?}
        An interface declares fields every implementor must have, so you
        can select common fields directly; a union's members share
        nothing, so all selection goes through inline fragments. Unions
        are the right tool for typed mutation results; interfaces for
        genuine ``is-a'' hierarchies.
  \item \textbf{What does the exclamation mark mean in \texttt{Int!}?}
        Non-null: the field always resolves to a value. The subtlety is
        the failure mode --- if a non-null field errors, the null bubbles
        to the nearest nullable ancestor rather than stopping at the
        field.
  \item \textbf{What are the three operation types?}
        Query (read), mutation (write; root fields execute serially, in
        order), subscription (selection plus an event stream, usually
        over WebSockets or SSE).
  \item \textbf{What are variables and why use them?}
        Typed arguments serialized in their own JSON object beside the
        query. They keep the document static and cacheable, and they make
        query-string injection structurally impossible --- there is no
        string context to escape from.
  \item \textbf{What is a fragment?}
        A named, reusable selection set. With \texttt{@include(if:)} and
        \texttt{@skip(if:)}, one cached operation can serve several UI
        states without string-templating the query.
  \item \textbf{What is an alias?}
        \texttt{active: products(status: ACTIVE)} --- renames the result
        key, letting you invoke the same field twice with different
        arguments in one operation.
  \item \textbf{What does a GraphQL response look like?}
        A JSON object with \texttt{data} (mirroring the operation,
        possibly partial or \texttt{null}) and, on field errors, an
        \texttt{errors} array whose entries carry \texttt{message},
        \texttt{path}, \texttt{locations}, and \texttt{extensions}.
        Field errors arrive as HTTP 200.
  % ---- mid ----
  \item \textbf{Explain nullability propagation on errors.}
        An error at a nullable field nulls that field. An error at a
        \emph{non-null} field cannot be represented --- the type forbids
        null --- so the error bubbles: the parent object is nulled, and
        so on up until a nullable position absorbs it; if none does,
        \texttt{data} itself becomes \texttt{null}. Non-null is a
        containment policy, not a reliability claim: it says ``if this
        fails, void my parent too.'' Design nullability from the client's
        rendering needs.
  \item \textbf{What is the N+1 problem in GraphQL and why is it
        structural?}
        Field-level resolution means every product's \texttt{category}
        resolver runs independently: a page of $N$ products plus one page
        query is $N{+}1$ store reads. It is structural because
        composability --- resolvers not knowing their caller --- is the
        feature that causes it. Our demo measures $1{+}2N$ reads naive
        versus 3 batched at $N{=}5$.
  \item \textbf{How does DataLoader solve N+1, and what is its caching
        scope?}
        Within one event-loop tick, resolvers enqueue keys instead of
        querying; the loader dedupes keys, issues one batch read, and
        distributes results positionally to the waiting futures. The
        cache is \textbf{per request}: one loader instance per request,
        built in the context factory. Sharing one across requests serves
        stale data across user boundaries --- a security bug, not a
        performance bug.
  \item \textbf{Describe the Relay connection spec.}
        A paginated field returns a connection: \texttt{edges}
        (opaque \texttt{cursor} + \texttt{node}), \texttt{pageInfo}
        (\texttt{hasNextPage}, \texttt{hasPreviousPage},
        \texttt{startCursor}, \texttt{endCursor}), with
        \texttt{first}/\texttt{after} forward and
        \texttt{last}/\texttt{before} backward. Cursors are value-based
        (keyset) bookmarks, computed with fetch-$p{+}1$ for
        \texttt{hasNextPage}; server caps \texttt{first}.
  \item \textbf{Field errors vs.\ top-level errors --- and where do typed
        mutation errors fit?}
        Top-level \texttt{errors} entries are exceptional failures
        (resolver could not produce the field). Expected business
        failures --- validation, rule rejection --- should be
        \emph{data}: a union payload of success and error types selected
        with inline fragments. Rule of thumb: typed errors inform the
        user; top-level errors page the on-call.
  \item \textbf{What goes in \texttt{extensions.code}?}
        The machine-readable error kind ---
        \texttt{COST\_DATA\_UNAVAILABLE}, \texttt{BAD\_CURSOR} --- so
        clients branch on codes, not on English \texttt{message} strings.
        It is the open extension point the spec deliberately leaves for
        exactly this.
  \item \textbf{How do you gate introspection, and why?}
        Reject documents selecting \texttt{\_\_schema} or
        \texttt{\_\_type} at the validation layer --- an AST check, not a
        substring match, because aliases and fragments defeat substrings.
        Why: introspection publishes your complete field inventory,
        including fields no client ships. Off for unauthenticated
        production traffic; on in development.
  \item \textbf{Why is mutation execution serial?}
        It is the specification's only ordering guarantee, made precisely
        so multi-step writes have defined semantics: root mutation fields
        run in document order, each seeing the previous one's effects.
        Queries have no such guarantee and may resolve in parallel.
  \item \textbf{What is a persisted query?}
        A pre-registered operation identified by hash. As an allow-list
        it closes the query language: unregistered documents never parse.
        As APQ it is also a bandwidth optimization --- send the hash,
        retry with full text only on \texttt{PersistedQueryNotFound}.
  \item \textbf{Walk through the APQ protocol flow.}
        Client sends \texttt{extensions.persistedQuery.sha256Hash} with
        no query text. Server misses: HTTP 200 with error code
        \texttt{PersistedQueryNotFound}. Client retries with query text
        plus hash; server verifies the hash matches the text
        (\texttt{PersistedQueryMismatch} if not), registers, executes.
        Steady state: 64 bytes uphill, and hash-keyed GETs become
        CDN-cacheable.
  % ---- senior ----
  \item \textbf{Design cost limiting for a public GraphQL API.}
        Layers, firing in order: (1) parse limits (document size);
        (2) static analysis between parse and execution --- depth cap,
        argument-weighted complexity (list fields multiply subtree cost
        by page size), alias and batch multipliers; (3) persisted-query
        allow-list for first-party clients; (4) per-field resolver
        timeouts as the backstop for estimation errors; (5) rate limiting
        keyed on \emph{complexity units per token}, not request count,
        because one request can carry a batch of fifty.
  \item \textbf{Why is request-count rate limiting insufficient for
        GraphQL?}
        The work per request varies by orders of magnitude --- the client
        composes it --- and array batching puts many operations in one
        HTTP request. ``100 requests/minute'' can mean 100 cheap reads or
        100 batches of god-queries. Limit estimated cost per principal
        per window; count requests only as a sanity backstop.
  \item \textbf{How do you version a GraphQL API?}
        Mostly, you don't: continuous evolution. Additive changes are
        free; remove nothing --- mark fields \texttt{@deprecated(reason:)}
        and track per-field usage telemetry until a field's callers reach
        zero, then remove. When a truly breaking redesign is unavoidable,
        version the \emph{endpoint} (\texttt{/graphql/v2}) as a last
        resort. The philosophy trades REST's flag days for an eternal
        compatibility obligation --- backed by telemetry, or it is a
        guess.
  \item \textbf{GraphQL subscriptions: transport choices?}
        WebSockets with the \texttt{graphql-ws} subprotocol is the
        mainstream: full-duplex, mature client support. Server-sent
        events are simpler (HTTP, auto-reconnect, CDN-friendly) and
        suffice for server-to-client streams. In WebSocket mode, remember
        authn/z at connection upgrade \emph{and} per-message, and cap
        concurrent subscriptions per principal --- each one holds server
        state.
  \item \textbf{Where does authorization live in a GraphQL server?}
        At the field, not the route --- the schema advertises everything
        to everyone, so ``the endpoint requires auth'' says nothing about
        \texttt{salaryHistory}. Enforce in resolvers or a directive /
        middleware layer with the principal on the context, deny by
        default, and audit the schema for fields lacking a declared
        scope. Where partial visibility is legitimate, return
        \texttt{NOT\_AUTHORIZED} as typed data.
  \item \textbf{``Variable injection is impossible'' --- true? What
        injection remains?}
        True for the query document: variables are a separate JSON
        object, never interpolated into the parse. Resolver-side
        injection is fully real: a resolver concatenating a variable into
        SQL, a shell command, or a downstream URL is injection in a typed
        costume. Parameterized queries and boundary validation remain
        non-negotiable \cite{kleppmann2017ddia}.
  \item \textbf{When would you refuse GraphQL and ship REST?}
        When the traffic is public, read-heavy, and CDN-cacheable ---
        GET with ETags is an operational superpower POST cannot match;
        when uploads/downloads dominate (multipart in GraphQL is an
        awkward bolt-on); when consumers are countless and unknown, so
        cost governance cannot rely on allow-lists; or when the team
        lacks the appetite for the guardrail stack GraphQL demands.
  \item \textbf{Your complexity estimator can be gamed. How, and what is
        the backstop?}
        Estimators price shapes, not data: a depth-3 query over a
        low-cardinality field can still fan out expensively in a
        resolver, and \texttt{first: 1} over a full-table scan is ``cheap''
        only on paper. Mitigations: per-field cost weights calibrated to
        measured resolver latency, timeouts per resolver and per
        operation, and query-count telemetry per operation name. Static
        analysis is the sieve; timeouts are the floor.
  % ---- staff ----
  \item \textbf{Design the error contract for a GraphQL platform serving
        thirty client teams.}
        Two documented channels: typed union payloads for expected,
        user-actionable failure (versioned like the schema, codes
        enumerated in the contract), and top-level errors with
        \texttt{extensions.code} for exceptional failure. A written
        nullability policy per domain (what voids the parent), client
        guidance on rendering partial data, and SLOs that count
        \texttt{errors} rates per code --- otherwise every team's alert
        channel becomes your schema review.
  \item \textbf{How do you roll out a schema change that is secretly
        breaking?}
        ``Breaking'' hides in semantics, not syntax: reinterpreting a
        field, tightening a non-null, reordering enum members a client
        switched on. Defense is process: additive-only CI checks on the
        SDL, per-field usage telemetry before any deprecation completes,
        contract tests generated from real persisted operations, and
        staged rollout with per-client pinning where the platform allows
        it.
  \item \textbf{When does federation pay for itself, and what new failure
        modes does it buy?}
        When schema ownership, not query volume, is the bottleneck ---
        many teams blocked on one schema's review queue. It buys:
        composition conflicts discovered at publish time, the gateway as
        a latency and availability chokepoint, entity resolution
        re-creating $N{+}1$ \emph{across services}, and distributed
        tracing becoming a prerequisite rather than a nicety. Federate
        for organizational scaling; not because the graph is big.
  \item \textbf{A client team demands \texttt{totalCount} on every
        connection. Take a position.}
        Refuse the default, offer the opt-in. \texttt{COUNT(*)} over a
        filtered relation is the offset problem in disguise --- cost
        grows with the corpus while the value is usually decorative
        (``1--50 of 38{,}412''). Ship \texttt{hasNextPage} by default,
        \texttt{totalCount} behind an explicit argument with a cost
        weight that prices it honestly, and approximate counts where the
        UI truly needs magnitude.
  \item \textbf{Design the migration of a 200-endpoint REST API to
        GraphQL without a flag day.}
        Strangler pattern: the GraphQL gateway fronts the existing REST
        services, resolvers initially delegate to them (and inherit their
        cacheability problems --- acknowledged, measured). Migrate
        consumer by consumer behind persisted operations; the allow-list
        doubles as the migration ledger. Deprecate REST endpoints only
        when telemetry shows zero traffic. The failure mode to design
        against is the permanent halfway state --- budget for finishing.
\end{enumerate}

%% ============================================================================
\section{External Bibliography \& Further Reading}
\label{sec:ch12-bibliography}

\begin{itemize}
  \item \textbf{GraphQL Specification (October 2021)} \cite{graphqlspec}
        --- the normative text. Read \S3 (type system), \S6 (execution),
        and \S7 (response) before trusting any tutorial, including this
        chapter.
  \item \textbf{graphql.org Learn} (graphql.org/learn) --- the
        foundation's tutorial track; the clearest prose introduction to
        schemas, queries, and mutations.
  \item \textbf{Relay GraphQL Cursor Connections Specification}
        (relay.dev/graphql/connections.htm) --- the normative
        \texttt{edges}/\texttt{pageInfo} contract our connection
        implements.
  \item \textbf{Apollo Server documentation: Persisted Queries and APQ}
        (apollographql.com/docs) --- the reference description of the
        hash-first handshake and its caching motivations.
  \item \textbf{Marc-Andr\'e Giroux, \emph{Production Ready GraphQL}}
        (2020) --- the practitioner's treatment of cost limiting,
        security, and schema design at scale; the closest book-length
        relative of this chapter.
  \item \textbf{Shopify's GraphQL API design documentation}
        (shopify.dev) --- a large public graph's design conventions;
        notable for connection discipline and cost-based rate limiting in
        production.
  \item \textbf{Kleppmann, \emph{Designing Data-Intensive Applications}}
        \cite{kleppmann2017ddia} --- the storage-side fundamentals that
        make the $N{+}1$ discussion precise.
  \item \textbf{RFC 9110, HTTP Semantics} \cite{rfc9110} --- status-code
        and method semantics for the GraphQL-over-HTTP mapping, and the
        caching machinery GraphQL opts out of by default.
  \item \textbf{OpenAPI Specification v3.1} \cite{openapi31} --- the
        contract document GraphQL renders unnecessary, useful for
        contrast in the REST-versus-GraphQL decision.
  \item \textbf{FastAPI and httpx documentation}
        \cite{fastapidocs,httpxdocs} --- the server host and client
        transport used in this chapter's listings.
\end{itemize}

%% ============================================================================
\section{Capstone Projects}
\label{sec:ch12-capstones}

Five projects, progressive. Each names its difficulty tier; acceptance
criteria are checkable facts, not vibes. All run offline against the
Northwind Robotics lab.

\subsection{Capstone 1: Catalog Schema from Scratch [junior]}

\textbf{Objective.} Re-implement the catalog schema without copying
Listing~\ref{lst:ch12-schema}: types, enum, input, union mutation result,
and a cursor connection, over your own seeded repository.

\textbf{Inputs/Datasets.} Synthetic catalog of your own design (10+
products, 3+ categories, reviews), deterministic seed.

\textbf{Steps.}
\begin{enumerate}
  \item Define \texttt{Product}, \texttt{Category}, \texttt{Review}, and
        the \texttt{ProductStatus} enum.
  \item Implement \texttt{products(first, after)} as a Relay connection
        with opaque decodable cursors and a \texttt{first} cap.
  \item Implement \texttt{addReview} returning a union of success and
        typed error payloads.
  \item Write in-process tests for the happy paths and one typed error.
\end{enumerate}

\textbf{Acceptance Criteria.}
\begin{itemize}
  \item[$\square$] A full cursor walk returns every product exactly once.
  \item[$\square$] A forged cursor yields a field error with a
        machine-readable \texttt{extensions.code}.
  \item[$\square$] Out-of-range rating returns the typed error payload
        with \texttt{errors} absent.
  \item[$\square$] The generated SDL (printed via the schema object)
        contains the union and the enum.
\end{itemize}

\textbf{Stretch Goals.} Add \texttt{@deprecated} to one field and verify
it appears in the introspected SDL; add backward pagination
(\texttt{last}/\texttt{before}).

\subsection{Capstone 2: N+1 Hunter [mid]}

\textbf{Objective.} Given a deliberately naive GraphQL service (naive
resolvers, no loaders), find and eliminate every $N{+}1$ with measured
proof.

\textbf{Inputs/Datasets.} Your Capstone-1 schema extended with
\texttt{Category.products} and \texttt{Review.author}; an instrumented
repository with a query counter.

\textbf{Steps.}
\begin{enumerate}
  \item Write three operations of increasing fan-out and record naive
        store-query counts.
  \item Introduce per-request DataLoaders for each relationship.
  \item Assert batched counts in tests --- the numbers are the
        deliverable.
  \item Demonstrate the per-request cache rule: a module-level loader
        serves stale data after a mutation; write the failing test, then
        fix by request-scoping.
\end{enumerate}

\textbf{Acceptance Criteria.}
\begin{itemize}
  \item[$\square$] Each operation's batched count is independent of page
        size (verified at two page sizes).
  \item[$\square$] A test proves duplicate keys are deduped into one
        batch read.
  \item[$\square$] The stale-cache regression test fails with a shared
        loader and passes with per-request loaders.
\end{itemize}

\textbf{Stretch Goals.} Add per-tick batch-size metrics and emit them in
a response \texttt{extensions} block.

\subsection{Capstone 3: Cost-Control Gateway [mid]}

\textbf{Objective.} Build the static analysis gate of
Listing~\ref{lst:ch12-cost} as a standalone, schema-agnostic library with
adversarial tests.

\textbf{Inputs/Datasets.} A corpus of operations: legitimate pages, deep
cyclic traversals, alias-multiplied god-queries, fragment-heavy
documents.

\textbf{Steps.}
\begin{enumerate}
  \item Implement depth and argument-weighted complexity over the parsed
        AST, expanding fragment spreads with cycle detection.
  \item Price aliases honestly: \texttt{a: products} and
        \texttt{b: products} both pay.
  \item Add a per-operation timeout at execution and prove it fires with
        a deliberately slow resolver.
  \item Emit rejection reasons as structured error extensions.
\end{enumerate}

\textbf{Acceptance Criteria.}
\begin{itemize}
  \item[$\square$] Each corpus operation is accepted or rejected
        exactly as labeled.
  \item[$\square$] Complexity is monotonically increasing in
        \texttt{first} (property-tested across sizes).
  \item[$\square$] Fragment cycles are rejected with a distinct code,
        never a hang or crash.
\end{itemize}

\textbf{Stretch Goals.} Per-field cost weights driven by measured
resolver latency; complexity reported to the client in response
extensions.

\subsection{Capstone 4: APQ Production Gateway [senior]}

\textbf{Objective.} Ship the lab as an operations-locked gateway: full
APQ, a strict allow-list mode, introspection policy, and a client that
survives every failure channel.

\textbf{Inputs/Datasets.} The chapter lab; a build-time operation
registry exported from a client repository (JSON hash $\to$ document).

\textbf{Steps.}
\begin{enumerate}
  \item Implement APQ both directions, including
        \texttt{PersistedQueryMismatch} defense.
  \item Add strict mode: only registered hashes execute, no query text
        accepted.
  \item Key rate limiting on summed complexity per principal per minute;
        reject array batches over a cap.
  \item Write the client: APQ with single retry, partial-data handling,
        and typed-error branching.
\end{enumerate}

\textbf{Acceptance Criteria.}
\begin{itemize}
  \item[$\square$] Strict mode rejects an unregistered document with a
        distinct code and zero resolver executions (proven by the query
        counter).
  \item[$\square$] The client completes the miss $\to$ register $\to$
        hash-only sequence and sends text only once.
  \item[$\square$] Complexity-based limiting admits small queries while
        rejecting a burst of god-queries under the same request count.
\end{itemize}

\textbf{Stretch Goals.} Serve registered read operations over GET for CDN
cacheability; document the cache-key design.

\subsection{Capstone 5: Federated Robotics Graph [staff]}

\textbf{Objective.} Split the catalog into two subgraphs --- catalog
(products, categories) and reputation (reviews) --- with
\texttt{Product} as a shared entity, and compose them behind a gateway
you implement in miniature, offline.

\textbf{Inputs/Datasets.} Two Strawberry schemas; a hand-rolled gateway
that parses operations, plans them across subgraphs, and stitches
results; the seeded datasets of the earlier capstones.

\textbf{Steps.}
\begin{enumerate}
  \item Define the entity boundary: catalog owns \texttt{Product},
        reputation extends it with \texttt{reviews}.
  \item Implement gateway query planning for the subset of the language
        your test operations need; say explicitly what is out of scope.
  \item Reproduce the \emph{cross-service} N+1 (entity resolution per
        product) and batch it.
  \item Write a composition check that fails CI when subgraphs disagree
        on the entity key.
\end{enumerate}

\textbf{Acceptance Criteria.}
\begin{itemize}
  \item[$\square$] One client operation spanning both subgraphs returns
        correctly stitched data.
  \item[$\square$] Cross-service entity resolution is batched; counts
        asserted in tests.
  \item[$\square$] A subgraph outage degrades only its fields (partial
        data, typed extensions), proven by a fault-injection test.
  \item[$\square$] A written scope note lists the federation features
        deliberately not implemented.
\end{itemize}

\textbf{Stretch Goals.} Add a subscription to the reputation subgraph and
carry it through the gateway; add per-subgraph cost budgets at the
gateway.

%% ============================================================================
\section{Python Dependencies Appendix}
\label{sec:ch12-dependencies}

\subsection{Pinned Requirements}

\begin{minted}{text}
# requirements.txt -- Chapter 12 GraphQL lab
strawberry-graphql[fastapi]>=0.230
fastapi>=0.110
httpx>=0.27
pytest>=8
\end{minted}

Install on Windows (PowerShell):

\begin{minted}{text}
python -m venv .venv
.\.venv\Scripts\Activate.ps1
pip install -r requirements.txt
pytest -q
\end{minted}

\subsection{strawberry-graphql[fastapi]}

\textbf{Description.} A code-first GraphQL library: Python dataclasses
and type hints define the schema; resolvers are plain (a)sync functions;
ships a DataLoader and framework integrations, FastAPI among them (the
\texttt{[fastapi]} extra).

\textbf{Why included.} Code-first keeps the chapter's listings teaching
the \emph{type system} rather than an SDL mapping layer; it is actively
maintained on modern \texttt{graphql-core}; its DataLoader and context
mechanics map directly onto the specification's execution model.

\textbf{Alternatives.} \textbf{Ariadne} --- SDL-first: you write the
schema in SDL and bind resolvers to it; choose it when the SDL must be
the reviewed artifact (multi-language organizations, contract-first
governance). \textbf{Graphene} --- the incumbent, similar code-first
style; broader legacy adoption, slower release cadence in recent years.
\textbf{GraphQL-core} alone --- the reference port; what Strawberry and
Graphene build on; too low-level to teach from directly.

\textbf{Installation (PowerShell).}
\begin{minted}{text}
pip install "strawberry-graphql[fastapi]>=0.230"
\end{minted}

\subsection{fastapi}

\textbf{Description.} The ASGI framework used to host the GraphQL
endpoint and health check \cite{fastapidocs}.

\textbf{Why included.} Consistency with Chapter~6, async-native request
handling, and a testable app object that \texttt{httpx.ASGITransport}
can execute in-process --- the offline backbone of every demo and test
in this chapter.

\textbf{Alternatives.} Starlette directly (FastAPI's base; fewer
conventions), or Strawberry's Django/Flask integrations where those
frameworks already own the service.

\textbf{Installation (PowerShell).}
\begin{minted}{text}
pip install "fastapi>=0.110"
\end{minted}

\subsection{httpx}

\textbf{Description.} A modern sync/async HTTP client with transport
injection \cite{httpxdocs}.

\textbf{Why included.} \texttt{httpx.ASGITransport} speaks real HTTP to
the in-process FastAPI app with zero sockets --- offline-first without
mocking semantics. The async client also carries the APQ handshake
cleanly.

\textbf{Alternatives.} \texttt{requests} (sync only, no ASGI transport),
\texttt{aiohttp} (async, different transport model), or a dedicated
GraphQL client such as \texttt{gql} (adds subscriptions support at the
price of another abstraction layer).

\textbf{Installation (PowerShell).}
\begin{minted}{text}
pip install "httpx>=0.27"
\end{minted}

\subsection{pytest}

\textbf{Description.} The test runner used across this book.

\textbf{Why included.} The suite's 21 tests run in-process and offline;
\texttt{asyncio.run} inside each test avoids any plugin dependency, so
\texttt{pytest} alone suffices.

\textbf{Alternatives.} \texttt{pytest-asyncio} (idiomatic async tests at
the price of a plugin), \texttt{unittest} (stdlib, more ceremony).

\textbf{Installation (PowerShell).}
\begin{minted}{text}
pip install "pytest>=8"
\end{minted}

%% ============================================================================
\section{Chapter Summary \& Key Takeaways}
\label{sec:ch12-summary}

\begin{itemize}
  \item GraphQL inverts control of response shape: the client composes,
        the server guarantees against a typed, introspectable schema. The
        same inversion makes the client choose the work --- cost
        governance is not optional.
  \item Nullability is error-containment policy expressed as syntax. An
        error at a non-null field bubbles to the nearest nullable
        position; the response is HTTP 200 with partial \texttt{data} and
        a populated \texttt{errors} array.
  \item Field-level resolution makes $N{+}1$ the default physics.
        DataLoader collapses it --- one loader per request, batched keyed
        reads, results aligned positionally. Our measured collapse: 11
        store reads to 3 at page size 5.
  \item Relay connections standardize cursor pagination:
        \texttt{edges}/\texttt{node}/\texttt{pageInfo}, opaque cursors,
        fetch-$p{+}1$ for \texttt{hasNextPage}, capped page sizes.
  \item Two error channels, chosen deliberately: typed union payloads for
        expected failure, top-level errors with \texttt{extensions.code}
        for exceptional failure.
  \item The guardrail stack, in firing order: depth limit, complexity
        budget, persisted-query allow-list (APQ for ergonomics and
        cacheable GETs), resolver timeouts, cost-keyed rate limiting,
        capped array batching, gated introspection, GraphiQL off in
        production.
  \item Variables make query-document injection impossible;
        resolver-side injection is unchanged from every other stack ---
        parameterized queries, always.
  \item REST keeps the crown for CDN-cacheable public reads, uploads, and
        unknown-consumer APIs; GraphQL wins for many-clients-one-backend
        product graphs; gRPC for internal low-latency RPC. Choose on
        constraints, not fashion.
\end{itemize}

%% ============================================================================
\section{Quick-Reference Card}
\label{sec:ch12-quickref}

\subsection{Type System at a Glance}

\begin{table}[h]
  \centering
  \small
  \begin{tabularx}{\textwidth}{l X}
    \toprule
    \textbf{Kind} & \textbf{Remember} \\
    \midrule
    Scalar & Leaves: \texttt{Int Float String Boolean ID} (+ custom) \\
    Object & Named field set; recursion makes the graph \\
    Interface & Shared field contract; inline fragments for the rest \\
    Union & No shared contract; \emph{all} selection via fragments \\
    Enum & Closed value set, validated pre-execution \\
    Input object & Argument-side only; no unions/interfaces inside \\
    List / Non-null & \texttt{[T]}, \texttt{T!}: error-containment policy \\
    Meta fields & \texttt{\_\_schema}, \texttt{\_\_type}, \texttt{\_\_typename} \\
    \bottomrule
  \end{tabularx}
  \caption{GraphQL type system cheat sheet.}
  \label{tab:ch12-typecard}
\end{table}

\subsection{Wire Shapes}

\begin{minted}{json}
// request
{ "query": "query P($first: Int!) { products(first: $first) { ... } }",
  "variables": { "first": 10 },
  "operationName": "P",
  "extensions": { "persistedQuery": { "version": 1, "sha256Hash": "..." } } }

// response: success, partial success, and pre-execution rejection
{ "data": { "products": { ... } } }
{ "data": { "product": null },
  "errors": [ { "message": "...", "path": ["product", "marginPercent"],
                "extensions": { "code": "COST_DATA_UNAVAILABLE" } } ] }
// HTTP 400 + errors array: parse / validation / budget rejection
\end{minted}

\subsection{Cost-Control Checklist}

\begin{itemize}
  \item[$\square$] Depth limit enforced on the parsed AST, pre-execution.
  \item[$\square$] Complexity budget with page-size multipliers; aliases
        and batch elements counted.
  \item[$\square$] Page sizes capped at the resolver (\texttt{first}
        $\le$ 50 here).
  \item[$\square$] Persisted-query allow-list for first-party clients;
        APQ for the rest.
  \item[$\square$] Array batching disabled or length-capped; cost summed
        across the batch.
  \item[$\square$] Rate limiting keyed on complexity units, not request
        count.
  \item[$\square$] Per-resolver and per-operation timeouts armed.
  \item[$\square$] Introspection gated (AST check), GraphiQL off in
        production.
  \item[$\square$] DataLoaders constructed per request, never shared.
  \item[$\square$] Field-level authorization deny-by-default, audited in
        CI.
\end{itemize}

\section{Standardized Evidence Addendum}
\subsection{Benchmark Snapshot Template}
\begin{table}[h]
  \centering
  \small
  \begin{tabularx}{\textwidth}{l c c c c}
    \toprule
    \textbf{Config} & \textbf{p50 latency} & \textbf{p99 latency} & \textbf{Error rate} & \textbf{Throughput} \\
    \midrule
    Baseline & -- & -- & -- & 1.00x \\
    Candidate A & -- & -- & -- & -- \\
    Candidate B & -- & -- & -- & -- \\
    \bottomrule
  \end{tabularx}
  \caption{Minimum benchmark table to keep per chapter.}
\end{table}

\subsection{Request Trace Template}
\begin{itemize}
  \item \textbf{Request:} one representative API call for this chapter's topic.
  \item \textbf{Before:} behavior (headers, payload, latency, failure mode) with the naive approach.
  \item \textbf{After:} behavior with the technique in this chapter.
  \item \textbf{Result:} what changed in correctness, resilience, latency, and operability.
\end{itemize}

\subsection{Cost and Latency Budget Snapshot}
\begin{table}[h]
  \centering
  \small
  \begin{tabularx}{\textwidth}{l c c}
    \toprule
    \textbf{Stage} & \textbf{Latency budget} & \textbf{Cost driver} \\
    \midrule
    TLS / connection setup & -- & handshake round trips, keep-alive hit rate \\
    Request validation & -- & schema complexity, CPU per request \\
    Business logic and I/O & -- & downstream fan-out, query cost \\
    Serialization and egress & -- & payload size, compression ratio \\
    \bottomrule
  \end{tabularx}
  \caption{Operational budget template for this chapter's technique.}
\end{table}

\section{Configuration Preset Template}
\begin{minted}{text}
# api_policy.yaml
policy_version: "v1.0.0"
component: "<chapter_topic>"
timeout_ms: 0
retry_budget: 0
fallback_strategy: "fail_fast"
rollback_trigger: "error_budget_burn_gt_threshold"
\end{minted}

\section{CI Quality Gate Specification}
\begin{itemize}
  \item contract tests green against the pinned OpenAPI/AsyncAPI/Protobuf spec,
  \item no breaking schema changes without a version bump,
  \item error responses conform to RFC 9457 problem-details shape,
  \item benchmark deltas within approved regression bounds (latency and throughput),
  \item policy version and rollback plan present in release metadata.
\end{itemize}

\section{Incident Runbook Template}
\begin{enumerate}
  \item Detect regression from SLO burn-rate alerts or consumer reports.
  \item Scope impacted endpoints, versions, and consumer segments.
  \item Contain impact (rollback, feature flag, rate-limit tightening, or traffic shift).
  \item Diagnose with traces, structured logs, and contract diffs.
  \item Recover with validated fix and verified rollout procedure.
  \item Prevent recurrence via new regression tests and CI gates.
\end{enumerate}

\section{Cross-Chapter Decision Card}
\begin{table}[h]
  \centering
  \small
  \begin{tabularx}{\textwidth}{l X X}
    \toprule
    \textbf{Dimension} & \textbf{Choose this approach when} & \textbf{Prefer an alternative when} \\
    \midrule
    Use case & -- & -- \\
    Latency budget & -- & -- \\
    Operational cost & -- & -- \\
    Team maturity & -- & -- \\
    \bottomrule
  \end{tabularx}
  \caption{Decision card to complete per chapter.}
\end{table}

\section{ROI Model and Build-vs-Buy}
\begin{itemize}
  \item \textbf{Build when:} the capability is differentiating, compliance forbids vendors, or scale makes vendor pricing irrational.
  \item \textbf{Buy when:} the capability is commodity (auth, gateways, observability), time-to-market dominates, or staffing is thin.
  \item \textbf{Hidden costs to model:} on-call load, upgrade cadence, expertise ramp, data egress fees.
\end{itemize}

\section{Disaster Recovery Notes (RPO/RTO)}
\begin{itemize}
  \item Define recovery point objective (RPO): maximum acceptable data loss window.
  \item Define recovery time objective (RTO): maximum acceptable restoration time.
  \item Test restores regularly; an untested backup is not a backup.
\end{itemize}

